\documentclass[11pt]{article}

\usepackage[letterpaper,margin=0.82in]{geometry}
\usepackage[T1]{fontenc}
\usepackage[utf8]{inputenc}
\usepackage{lmodern}
\usepackage{microtype}
\usepackage{amsmath,amssymb,amsthm,mathtools}
\usepackage{booktabs,longtable,array}
\usepackage{enumitem}
\usepackage{xcolor}
\usepackage{hyperref}
\usepackage{fancyhdr}
\usepackage{listings}
\usepackage{graphicx}
\usepackage{float}
\usepackage{caption}
\captionsetup{font=small,labelfont=bf}
\graphicspath{{figures_final/}}

\hypersetup{
  colorlinks=true,
  linkcolor=blue!45!black,
  citecolor=blue!45!black,
  urlcolor=blue!45!black,
  pdftitle={Komorebi Quadratic-Form Engine and K-COS: Final Traditional Edition},
  pdfauthor={Komorebi P.}
}

\setlength{\parindent}{0pt}
\setlength{\parskip}{6pt}
\setlength{\headheight}{14pt}
\setlength{\emergencystretch}{2em}
\makeatletter
\renewcommand*\l@subsection{\@dottedtocline{2}{1.5em}{3.1em}}
\renewcommand*\l@subsubsection{\@dottedtocline{3}{4.6em}{4.2em}}
\makeatother

\newtheorem{theorem}{Theorem}[section]
\newtheorem{proposition}[theorem]{Proposition}
\newtheorem{corollary}[theorem]{Corollary}
\newtheorem{lemma}[theorem]{Lemma}
\theoremstyle{definition}
\newtheorem{definition}[theorem]{Definition}
\newtheorem{algorithm}[theorem]{Algorithm}
\theoremstyle{remark}
\newtheorem{remark}[theorem]{Remark}
\DeclareMathOperator{\E}{\mathbb E}
\DeclareMathOperator{\Var}{Var}
\newcommand{\R}{\mathbb R}
\newcommand{\C}{\mathbb C}
\newcommand{\Pp}{\mathbb P}
\newcommand{\eps}{\varepsilon}
\newcommand{\ii}{\mathrm{i}}
\newcommand{\phit}{\phi_{\theta}}


\pagestyle{fancy}
\fancyhf{}
\fancyhead[L]{Komorebi Quadratic-Form Engine}
\fancyhead[R]{Komorebi P.}
\fancyfoot[C]{\thepage}

\title{\textbf{Komorebi Quadratic-Form Engine and K-COS}\\
\large Memory-Bounded Structured Solvers, Query-Adaptive Spectral Packets,\\
and Directed-Rounding Extreme-Tail Enclosures}
\author{Komorebi P.}
\date{Final traditional edition --- October 7, 2026}


% Operation Komorebi: shared publication profile.
\usepackage{komorebi-publication}
\renewcommand{\KomorebiTitle}{Quadratic-Form Engine and K-COS}
\renewcommand{\KomorebiSubtitle}{Memory-bounded structured solvers, query-adaptive spectral packets, and directed-rounding extreme-tail enclosures}
\renewcommand{\KomorebiSeries}{Research library}
\renewcommand{\KomorebiKind}{Traditional paper}
\renewcommand{\KomorebiEdition}{Standardized publication edition / 8 October 2026}
\renewcommand{\KomorebiShortTitle}{Quadratic-Form Engine and K-COS}
\renewcommand{\KomorebiTitleSize}{38}
\renewcommand{\KomorebiTitleLeading}{43}
\renewcommand{\KomorebiProfile}{classical}
\renewcommand{\KomorebiOrientation}{portrait}
\KomorebiSetup
\begin{document}
\KomorebiCover

\relax

\begin{abstract}
Gaussian quadratic forms occur throughout statistics, numerical analysis,
signal processing, kernel methods, and variance-component inference.  This
work develops a single computational framework for central positive,
central indefinite, and positive noncentral forms, with particular emphasis
on two difficult regimes: operators too large to reside in accelerator memory
and survival probabilities many hundreds or thousands of decades into the
tail.

The method separates storage, structural propagation, rank compression, and
spectral tail evaluation.  Exact block identities permit out-of-core factor
actions without changing the represented operator.  Polyphase light-cone
theorems and certified approximate light cones turn structural locality into
phase-aware scheduling.  Range-residual invariants quantify unresolved
spectral mass and trigger exact residual completion when the remaining rank is
small.  The final tail layer uses exponential tilting, query-adaptive spectral
packets, log-domain accumulation, and explicit support and Fourier-mode
remainder bounds.

For exact supplied positive spectra, a directed-rounding backend returns a
logarithmic probability interval and an accuracy-driven controller refines
internal resolution until the requested multiplicative probability accuracy is
proved or fails closed.  In the archived September 2026 exact control
\(0.02\chi^2_{80}\), the solver reported tail enclosures through \(10^{-5000}\);
at the \(10^{-5000}\) point a \(10^{-18}\) requested relative accuracy is
met with a certified ratio-minus-one of approximately
\(9.2\times10^{-20}\).  A fresh October 7, 2026 run confirms the deepest control with directed
ratio upper endpoint below \(9.202\times10^{-20}\); an independent exact
Erlang interval is fully contained.  The final regression suites pass 58 of
58 tests.  Additional archived validation includes a rank-1997
production-scale positive spectrum near \(10^{-922}\), a corrected
strong-dependence spectrum near \(10^{-458}\), signed finite-endpoint controls
through \(10^{-200}\), a rank-340 mixed-sign tail near \(10^{-82}\), and
noncentral exact controls through \(10^{-40}\).

The central design rule is that approximation is an internal computational
choice, not the accuracy contract.  A machine-enclosed result is accepted only when its directed-rounding
probability interval proves the requested accuracy.  A binary64 analytic
ledger bounds mathematical approximation error and remains a distinct
evidence level because it does not enclose numerical rounding.
\end{abstract}


\section*{Final-edition contract and reading map}
\addcontentsline{toc}{section}{Final-edition contract and reading map}
This edition preserves the full matrix-to-tail manuscript from the v1.9
source archive and integrates the complete centered-packet mathematics from
the earlier K-COS paper.  The scope is a domain-neutral quadratic-form
engine; the rare-variant case study is one application.

Three companion objects accompany this traditional text: an explanatory
visual guide, executable Python source and fresh validation, and a focused
Lean companion.  The Lean companion contains 26 freshly kernel-checked finite algebraic and
certificate-composition theorem declarations; it does not formally certify the entire
probability analysis or the numerical libraries.

Unless a result is explicitly dated October 7, 2026, all numerical
measurements, including quantitative prose claims, tables, and plots in the
main body are \emph{archived September 2026 development evidence}.
They remain intact to preserve the research record.  Their precision,
implementation versions, and reference convergence are those documented in
the original source archive.  Fresh checks are separately identified in
Appendix~\ref{app:fresh-validation}; they do not silently replace historical
measurements or establish an unmeasured hardware speedup.

The word ``certificate'' has three distinct meanings here.  Structural
certificates are exact inequalities proved in the real-arithmetic model.
Analytic probability ledgers bound packet, support, and Fourier truncation
errors in that model.  Directed-rounding probability enclosures additionally
bound all finite numerical arithmetic in the specific interval backend.
Binary64 signed, noncentral, and packet calculations are analytic-ledger
results, even when their historical API returned \texttt{certified=True}.
That Boolean alone is not a floating-point enclosure.

The directed positive backend targets an \emph{exactly supplied} positive
spectrum of rank greater than two, a nonnegative threshold, a valid positive
tilt, a strictly larger positive support endpoint, and at least one retained
cosine mode.  The accuracy-driven designer currently uses upper-tail queries
above the spectral mean.  Matrix input certification targets the exact Gram
matrix of the supplied binary64 entries.  It does not certify unrecorded
preprocessing, a statistical model fit, or a real-valued matrix preceding the
stored entries.

\begin{center}
\small
\begin{tabular}{@{}p{.23\linewidth}p{.33\linewidth}p{.35\linewidth}@{}}
\toprule
Reading path & Mathematical content & Companion and evidence\\
\midrule
Representation and structure & Exact factor identities, storage refinement,
phase propagation, QR/range cores & Full proofs here; finite algebraic subset
in Lean; original numerical invariance checks\\[2mm]
Spectral uncertainty & Norm/Weyl bounds, residual moments, exact completion,
dense eigensystem post-verification & Proofs here; moderate-dimension directed
backend; scalable verification remains open\\[2mm]
K-COS packets & Centered complex CGF expansion, query grid, folding,
packet/support/mode error budget & Complete packet theory in
Section~\ref{sec:packets}; analytic and binary64 evidence kept separate\\[2mm]
Probability accuracy & Positive directed enclosure and tolerance acceptance &
Interval-library correctness is a software dependency; fresh regression
checks and archived extreme-tail controls\\
\bottomrule
\end{tabular}
\end{center}

\tableofcontents
\clearpage

\section{Results at a glance}\label{sec:final-1}

The engine is organized around a simple contract: preserve the quadratic form
while changing its representation, and report a tail probability only at an
accuracy that has been explicitly established.

\begin{center}
\small
\begin{tabular}{@{}p{0.27\linewidth}p{0.64\linewidth}@{}}
\toprule
Component & Resolved statement\\
\midrule
Central positive form &
\(X^\top M X\), with \(X\sim N(0,\Sigma)\) and \(M\succeq0\), reduces
exactly to \(\sum_j\lambda_j Z_j^2\).\\[2mm]

Central indefinite form &
For symmetric \(M\), the same reduction holds with signed
\(\lambda_j\), yielding a signed generalized chi-square form.\\[2mm]

Positive noncentral form &
\(X\sim N(\mu,\Sigma)\), \(M\succeq0\), reduces exactly to a deterministic
shift plus \(\sum_j\lambda_j\chi^2_1(\delta_j)\).\\[2mm]

Memory-bounded factorization &
At fixed resident column width \(c\), nuisance-space rank \(p\), and Krylov
width \(s\), streamed sample-space actions use
\(O(nc+ns+cs+np)\) working storage rather than \(O(nm)\).\\[2mm]

Structural scheduling &
Exact and approximate phase light cones certify which operator blocks can
affect a finite-depth Krylov query.\\[2mm]

Spectral recovery &
Range residual trace, square-trace, and spectral-cap information bounds
unresolved spectrum; when the remaining rank is small, streamed residual
completion restores the exact nonzero spectrum.\\[2mm]

Extreme-tail contract &
For exact positive spectra, the public input is a requested multiplicative
accuracy \(\varepsilon\).  The solver returns
\([p_-,p_+]\) only as certified when
\(p_+/p_- - 1\le\varepsilon\).\\[2mm]

Deepest rigorous control &
For \(0.02\chi^2_{80}\) at \(p=10^{-5000}\), a requested
\(10^{-18}\) relative accuracy is certified with achieved
ratio-minus-one approximately \(9.2\times10^{-20}\).\\
\bottomrule
\end{tabular}
\end{center}

The remainder of the manuscript develops these statements from the operator
level down to the final probability enclosure.


\section{General quadratic-form core}\label{sec:gaussian-core}

The primary mathematical object is a Gaussian quadratic form.

\subsection{Central positive forms}

Let
\[
X\sim N(0,\Sigma),
\qquad
M\succeq0,
\]
and write \(S=\Sigma^{1/2}\).  With \(X=SZ\), \(Z\sim N(0,I)\),
\[
Q=X^\top MX
=
Z^\top(SMS)Z.
\]
If the nonzero eigenvalues of the positive semidefinite operator
\[
H=SMS
\]
are \(\lambda_1,\ldots,\lambda_r\), then
\[
\boxed{
Q\overset d=
\sum_{j=1}^r\lambda_jZ_j^2.
}
\]

The numerical tail problem therefore depends only on a positive spectrum,
regardless of how that spectrum was generated.

\begin{proposition}[Representation equivalence]\label{claim:final-2-1}
The same central quadratic-form law may be supplied by any of the following:
\begin{enumerate}
\item a nonnegative spectrum \(\{\lambda_j\}\);
\item a positive semidefinite matrix \(H\);
\item a factor \(A\) with \(H=AA^\top\) or \(H=A^\top A\);
\item a replayable streamed factor
\[
A=[A_1\ \cdots\ A_B].
\]
\end{enumerate}
The nonzero spectra of \(AA^\top\) and \(A^\top A\) are identical, and
\[
AA^\top=\sum_bA_bA_b^\top.
\]
\end{proposition}

This proposition is the domain-independent source of the memory-bounded
streaming method developed below.

\subsection{Central indefinite forms}

The positive case is not the natural boundary of the central theory.  Let
\[
X\sim N(0,\Sigma),
\qquad
M=M^\top,
\]
where $M$ may be indefinite, and again write $S=\Sigma^{1/2}$.  The symmetric
operator
\[
H=SMS
\]
has a real signed spectrum.

\begin{theorem}[Exact central signed reduction]\label{claim:final-2-2}
If $\lambda_1,\ldots,\lambda_r$ are the nonzero eigenvalues of $H$, then
\[
\boxed{
X^\top M X
\overset d=
\sum_{j=1}^r\lambda_jZ_j^2.
}
\]
\end{theorem}

\begin{proof}
Write $X=SZ$ with $Z\sim N(0,I)$ and diagonalize
\[
SMS=V\Lambda V^\top.
\]
Rotational invariance of the standard normal gives
\[
Z^\top V\Lambda V^\top Z
\overset d=
\sum_j\lambda_jZ_j^2.
\]
No positivity assumption on $M$ is needed for this central reduction.
\end{proof}

The sign pattern determines the support.  If every $\lambda_j\ge0$, the lower
endpoint is zero.  If every $\lambda_j\le0$, the upper endpoint is zero.  If
both signs occur, the distribution is unbounded in both directions.

For the right-tail path, the cumulant generating function is
\[
K(s)
=
-\frac12\sum_j\log(1-2s\lambda_j),
\]
with real domain
\[
\boxed{
\max_{\lambda_j<0}\frac{1}{2\lambda_j}
< s <
\min_{\lambda_j>0}\frac{1}{2\lambda_j}.
}
\]
Here a missing negative or positive spectral set gives a domain endpoint
\(-\infty\) or \(+\infty\), respectively.
When $q>E[Q]=\sum_j\lambda_j$, the saddlepoint is the unique positive solution
of
\[
K'(\theta)=q
\]
inside this domain whenever a positive spectral component is present.  The
finite-upper-endpoint negative-only case is handled by the same tilted
geometry with the endpoint fixed at zero.

The packet representation is sign aware: positive and negative eigenvalues are
clustered separately, and each same-sign packet stores a center $\mu$ and
centered moments.  The local expansion remains
\[
-\frac12\log(1-2s(\mu+\delta))
=
-\frac12\log(1-2s\mu)
+
\frac12\sum_{n\ge1}
\frac{\alpha(s)^n\delta^n}{n},
\qquad
\alpha(s)=\frac{2s}{1-2s\mu},
\]
with the same geometric remainder condition
\[
|\alpha(s)|\max|\delta|<1.
\]

A useful distinction appears in the Fourier tail.  Under the saddlepoint tilt,
\[
\left|\phi_\theta(t)\right|
=
\prod_j
\left[
1+
\left(
\frac{2\lambda_j t}{1-2\theta\lambda_j}
\right)^2
\right]^{-1/4}.
\]
If $r$ coefficients are nonzero, then asymptotically
\[
|\phi_\theta(t)|=O(|t|^{-r/2}).
\]
Thus certification difficulty depends strongly on spectral rank, not only on
how small the requested probability is.

\paragraph{Exact finite-endpoint control.}
For
\[
Q=-0.02\,\chi^2_{80},
\]
the right tail is a lower chi-square tail with finite upper endpoint zero.  The
signed engine agrees with the exact distribution through
\[
\boxed{p=10^{-200}},
\]
with maximum observed log-probability error about $2.84\times10^{-14}$ decades
and certificates around $10^{-12}$.

\paragraph{Mixed-sign high-rank analytic-ledger control.}
A deterministic rank-340 multiscale spectrum with 180 positive and 160
negative coefficients was tested from approximately $10^{-14}$ to
\[
\boxed{p\approx10^{-81.891937}}.
\]
At the deepest point the packetized solver used 35 packets, $K=4096$, and
reported a relative certificate
\[
5.78\times10^{-9}.
\]
Its log-probability differed from the unpacketized full-spectrum signed solver
by about $1.42\times10^{-12}$ decades.  Independent exponential-tilt
importance sampling placed the discrepancy at about $1.41$ Monte Carlo
standard errors.

\paragraph{Retained low-rank limitation.}
The exact mixed-sign identity
\[
Q=0.07(\chi^2_2-\chi^2_2)
\]
is Laplace distributed with scale $0.14$, so for $q\ge0$,
\[
P(Q\ge q)=\frac12e^{-q/0.14}.
\]
The computed values agree with this exact law through $10^{-200}$ to roughly
$10^{-9}$ decades or better.  However, the current analytic COS-series
certificate rejects these low-rank mixed cases because the tilted
characteristic function decays only as $|t|^{-2}$.  The implementation reports
this as a certificate failure rather than silently accepting the numerical
value.

\begin{figure}[H]
\centering
\includegraphics[width=.84\linewidth]{figure_24_signed_certificate_boundary.png}
\caption{Signed-form certification boundary.  Negative finite-endpoint forms
and the rank-340 mixed spectrum remain certifiable in deep tails.  The exact
rank-4 Laplace control remains numerically accurate but fails the current
analytic COS-series certificate because its Fourier decay is much slower.}
\label{fig:signed-certificate-boundary}
\end{figure}

\subsection{Noncentral positive forms}

Now let
\[
X\sim N(\mu,\Sigma),
\qquad
M\succeq0.
\]
Again set \(S=\Sigma^{1/2}\) and diagonalize
\[
SMS=V\Lambda V^\top.
\]
Writing
\[
c=V^\top SM\mu,
\]
each positive-eigenvalue direction satisfies
\[
\lambda_iY_i^2+2c_iY_i
=
\lambda_i
\left(
Y_i+\frac{c_i}{\lambda_i}
\right)^2
-
\frac{c_i^2}{\lambda_i}.
\]

\begin{theorem}[Exact noncentral reduction]\label{claim:final-2-3}
For the positive eigenvalues \(\lambda_i>0\),
\[
\boxed{
Q
\overset d=
q_0
+
\sum_i
\lambda_i
\chi^2_1(\delta_i),
}
\]
where
\[
\delta_i=
\left(\frac{c_i}{\lambda_i}\right)^2
\]
and
\[
q_0
=
\mu^\top M\mu
-
\sum_i\frac{c_i^2}{\lambda_i}
\ge0.
\]
If \(\Sigma\) is nonsingular on the \(M\)-active subspace, then \(q_0=0\).
When \(\Sigma\) is singular, \(q_0\) is the deterministic quadratic
contribution outside the random covariance support.
\end{theorem}

\begin{proof}
Expand
\[
(\mu+SZ)^\top M(\mu+SZ)
=
Z^\top SMSZ
+
2\mu^\top MSZ
+
\mu^\top M\mu
\]
and rotate by \(V\).  Completing the square gives the displayed
noncentral components.  For a null vector of \(SMS\),
\[
M^{1/2}Sv=0,
\]
hence
\[
v^\top SM\mu
=
(M^{1/2}Sv)^\top M^{1/2}\mu
=
0,
\]
so there is no unresolved stochastic linear term in a zero-eigenvalue
direction.  The remaining constant is the squared norm of the deterministic
component and is therefore nonnegative.
\end{proof}

The production API implements this reduction directly.  In an isotropic
control,
\[
\Sigma=I,\qquad M=aI,
\]
the law collapses to the exact scaled noncentral chi-square distribution
\[
Q/a\sim\chi^2_d(\|\mu\|^2).
\]

\begin{figure}[H]
\centering
\includegraphics[width=.82\linewidth]{figure_22_noncentral_controls.png}
\caption{Independent noncentral controls.  Three scaled isotropic Gaussian
quadratic forms with different dimensions and noncentralities were compared
against SciPy's noncentral-\(\chi^2\) law from \(10^{-4}\) through
\(10^{-40}\).  The largest observed discrepancy was approximately
\(9.3\times10^{-15}\) decades, with all analytic truncation certificates
passing.}
\label{fig:noncentral-controls}
\end{figure}

\subsection{Decomposition axes}

The solver separates four logically different representations:
\[
\boxed{
\text{storage blocks}
\circ
\text{structural phases}
\circ
\text{Krylov/range compression}
\circ
\text{spectral tail packets}.
}
\]
Storage blocks answer where data reside.  Structural phases answer how the
operator propagates.  Krylov/range compression answers which subspace has
actually been resolved.  Spectral packets answer how the final transform is
evaluated.  These decompositions should not be conflated.


\section{Demonstrated deep-tail precision}\label{sec:final-3}
\textbf{Evidence date and scope.} All measurements in this section are
archived September 2026 runs.  Only rows explicitly described as directed
rounding are numerical enclosures.  The other reported certificates are
analytic approximation ledgers evaluated in floating point.

The method is designed for a regime in which direct survival probabilities can
underflow long before the mathematical probability becomes numerically
unresolvable.  The engine therefore treats \(\log p\) as the primary numerical
quantity and keeps approximation certificates alive in the same extreme-tail
regime.

The validation ladder spans seven qualitatively different settings:

\begin{center}
\small
\begin{tabular}{@{}p{0.27\linewidth}p{0.13\linewidth}p{0.22\linewidth}p{0.16\linewidth}p{0.12\linewidth}@{}}
\toprule
validation regime & deepest $\log_{10}p$ & observed error & certificate & spectral size\\
\midrule
directed-rounding $0.02\chi^2_{80}$ control
& $-5000$ & enclosure width $3.37\times10^{-11}$ decades
& multiplicative ratio $\le 1+7.75\times10^{-11}$ & 80 df\\
production-scale positive spectrum
& $-922.031162$ & below stored reference resolution & $5.71\times10^{-9}$ & rank 1997\\
corrected strong-dependence case study
& $-457.973916$ & below stored reference resolution & $5.20\times10^{-9}$ & rank 996\\
heterogeneous inversion benchmark near $10^{-16}$
& $-16$ & $2.88$--$6.52\times10^{-15}$ decades & independent integral converged & ranks 80--120\\
signed finite-endpoint exact control
& $-200$ & $\le 2.84\times10^{-14}$ decades & $7.23\times10^{-13}$ & rank 80\\
signed mixed multiscale
& $-81.891937$ & $1.42\times10^{-12}$ decades vs full grid & $5.78\times10^{-9}$ & rank 340\\
noncentral exact controls
& $-40$ & $\le 9.26\times10^{-15}$ decades & passing & up to 50 df\\
\bottomrule
\end{tabular}
\end{center}

The analytic scaled-chi-square ladder also includes
\[
\log_{10}p=-100,-300,-1000,-5000,
\]
with respective reported absolute log-probability discrepancies
\[
2.84\times10^{-14},\quad
<10^{-15},\quad
3.41\times10^{-13},\quad
<10^{-15}
\]
in the stored benchmark output.  The corresponding analytic relative certificates
remain finite throughout, reaching only
\[
3.85\times10^{-11}
\]
at \(10^{-5000}\).

A production-scale run with \(n=2000\), \(m=50{,}000\), 49,951 retained
features, and quadratic-form rank 1997 produced
\[
\boxed{\log_{10}p=-922.0311624669905}
\]
using 79 spectral packets and \(K=128\).  Its independently evaluated
full-spectrum tilted reference agreed within the stored reporting precision, while the
reported relative certificate was
\[
5.71\times10^{-9}.
\]

The corrected strong-dependence case study independently reached
\[
\boxed{\log_{10}p=-457.9739160722001}
\]
with a finite \(5.20\times10^{-9}\) certificate.  The signed finite-endpoint
control passed its analytic ledger through \(10^{-200}\), while a rank-340 mixed-sign
multiscale spectrum passed its analytic ledger to
\[
\boxed{\log_{10}p=-81.89193686056221}.
\]
Finally, the noncentral engine matched exact noncentral-chi-square controls
from \(10^{-4}\) through \(10^{-40}\), with maximum observed discrepancy
\[
9.26\times10^{-15}
\]
decades.

\begin{figure}[H]
\centering
\includegraphics[width=.86\linewidth]{figure_23_deep_tail_precision_envelope.png}
\caption{Demonstrated precision envelope across positive analytic controls,
a production-scale rank-1997 spectrum, corrected strong-dependence spectra,
signed finite-endpoint and mixed-sign forms, and noncentral exact controls.  Points whose stored benchmark difference rounded to zero are
placed at $10^{-16}$ decades only to make them visible on the logarithmic
vertical axis.}
\label{fig:deep-tail-precision-envelope}
\end{figure}

\subsection{Independent characteristic-function inversion near \texorpdfstring{\(10^{-16}\)}{1e-16}}

Closed-form controls are important, but they can hide numerical difficulties
that appear only for heterogeneous spectra.  An independent inversion test was
therefore run on four generalized chi-square laws near
\[
P(Q\ge q)\approx10^{-16}.
\]

The reference evaluates the original characteristic function directly through
the central Imhof inversion identity
\[
P(Q\ge q)
=
\frac12
+
\frac1\pi
\int_0^\infty
\frac{\sin\Theta(u)}{u\,R(u)}\,du,
\]
where, for distinct positive coefficients \(\lambda_j\) with multiplicities
\(h_j\),
\[
\Theta(u)
=
\frac12\sum_j h_j\arctan(\lambda_j u)
-
\frac12qu,
\qquad
R(u)
=
\prod_j
\left(1+\lambda_j^2u^2\right)^{h_j/4}.
\]
This route does not use exponential tilting, the COS expansion, spectral
packets, or the K-COS error controller.

The four spectra include balanced, skewed, dominant-spike, and higher-rank
cases.  Each reference integral was evaluated at two truncation limits.  The
omitted integral after \(U\) was bounded using
\[
\frac1\pi
\int_U^\infty\frac{du}{uR(u)}
\le
\frac{C}{\pi p}U^{-p},
\qquad
p=\frac12\sum_jh_j,
\qquad
C=\prod_j\lambda_j^{-h_j/2}.
\]

\begin{center}
\small
\begin{tabular}{@{}lrrrr@{}}
\toprule
spectrum & rank & \(|\Delta\log_{10}p|\) & relative difference & reference tail bound / \(p\)\\
\midrule
balanced & 80 & $6.52\times10^{-15}$ & $1.50\times10^{-14}$ & $6.84\times10^{-20}$\\
skewed & 80 & $4.13\times10^{-15}$ & $9.50\times10^{-15}$ & $4.44\times10^{-18}$\\
spiky & 80 & $5.19\times10^{-15}$ & $1.20\times10^{-14}$ & $7.00\times10^{-17}$\\
balanced & 120 & $2.88\times10^{-15}$ & $6.64\times10^{-15}$ & $4.19\times10^{-29}$\\
\bottomrule
\end{tabular}
\end{center}

The two inversion truncations agree to between approximately
\[
1.2\times10^{-21}
\quad\text{and}\quad
3.7\times10^{-17}
\]
decades across the four cases.  Thus the slow reference is numerically stable
at a scale well below the observed K-COS-versus-inversion difference.

\begin{figure}[H]
\centering
\includegraphics[width=.82\linewidth]{figure_27_imhof_p1e16_validation.png}
\caption{Independent direct-inversion validation near $p=10^{-16}$ on four
heterogeneous spectra.  Circles show the relative K-COS/inversion difference;
crosses show the analytic bound on the inversion integral omitted beyond the
final truncation point.}
\label{fig:imhof-p1e16}
\end{figure}

This test is complementary to the exact scaled-chi-square controls.  It asks
whether the solver remains accurate when no single-gamma reduction is
available and when an algorithmically different inversion of the full
characteristic function is used as the reference.

These results have different evidentiary roles.  The scaled-chi-square,
finite-endpoint signed, Laplace signed, and noncentral cases provide analytic
controls.  The heterogeneous inversion battery, production-scale, strong-dependence, and
rank-340 mixed-sign cases test irregular spectra for which no single-gamma
closed form is available.  Together they test correctness against known distributions,
representation compression, and stability on the intended high-rank numerical
regime.  The low-rank Laplace example also records a deliberate certificate
failure boundary rather than conflating numerical agreement with certification.

\section{Directed-rounding deep-tail enclosure}\label{sec:final-4}

The preceding precision experiments combine analytic approximation bounds with
stable floating-point evaluation.  A stronger numerical statement is possible:
return an interval that is guaranteed to contain the probability represented
by an exact supplied positive spectrum.

This section gives the first such backend.

\subsection{Interval target}

For a positive central quadratic form
\[
Q=\sum_{j=1}^{r}\lambda_jZ_j^2,
\qquad
\lambda_j>0,
\]
the output is not merely a point estimate.  It is
\[
\boxed{
\log_{10}P(Q\ge q)\in[d_-,d_+].
}
\]
Hence the corresponding multiplicative uncertainty satisfies
\[
\boxed{
\frac{p_+}{p_-}\le 10^{d_+-d_-}.
}
\]

The finite transform and exponential prefactor are evaluated by
directed-rounding interval arithmetic.  The resulting interval is enlarged by
two independently proved analytic remainders: a support/periodization alias
bound and an omitted-COS-mode bound.

\subsection{No exact saddlepoint is required}

Choose any
\[
0<\theta<\frac{1}{2\lambda_{\max}}.
\]
Exponential tilting gives
\[
P(Q\ge q)
=
e^{K(\theta)-\theta q}
S_\theta(q),
\]
where
\[
S_\theta(q)
=
E_\theta\!\left[
e^{-\theta(Q-q)}
1_{\{Q\ge q\}}
\right].
\]

The interval proof does not require \(\theta\) to be the exact saddlepoint.
A high-precision saddlepoint is used only as an efficient design choice.
Once serialized as an exact decimal input, it is simply a valid tilt inside
the MGF domain.

The same is true of the support endpoint \(U\) and the auxiliary Chernoff
parameter \(\eta\).  Correctness requires only
\[
U>q\ge0,\qquad U>0,
\qquad
\theta<\eta<\frac{1}{2\lambda_{\max}}.
\]

\subsection{Rigorous support and alias control}

Under the tilted measure,
\[
P_\theta(Q>U)
\le
\exp\left\{
K(\eta)-K(\theta)-(\eta-\theta)U
\right\}.
\]

The full-characteristic-function cosine expansion periodizes the tilted density
on \([0,U]\).  Because the original positive quadratic form is supported on
\([0,\infty)\), the total positive-tail mass aliased back into one fundamental
interval is at most \(P_\theta(Q>U)\).

The actual payoff omitted above \(U\) is at most
\[
e^{-\theta(U-q)}P_\theta(Q>U).
\]

Therefore
\[
\boxed{
E_{\rm support}
\le
P_\theta(Q>U)
\left(
1+e^{-\theta(U-q)}
\right).
}
\]

\subsection{Rigorous infinite-mode control}

Define
\[
a_j=
\frac{2\lambda_j}{1-2\theta\lambda_j}.
\]
The normalized tilted characteristic function satisfies
\[
|\phi_\theta(t)|
=
\prod_{j=1}^r
\left(
1+(a_jt)^2
\right)^{-1/4}.
\]

Since
\[
1+(a_jt)^2\ge(a_jt)^2,
\]
we have
\[
|\phi_\theta(t)|
\le
C_\theta t^{-r/2},
\qquad
C_\theta=\prod_j a_j^{-1/2}.
\]

For the payoff coefficient
\[
I_k=
\int_q^U
e^{-\theta(x-q)}
\cos(k\pi x/U)\,dx,
\]
the elementary bound
\[
|I_k|\le\frac1\theta
\]
holds.

Let
\[
p=\frac r2>1,\qquad K\ge1.
\]
For a finite transform retaining modes \(0,\ldots,K-1\),
\[
\boxed{
E_{\rm series}
\le
\frac{2C_\theta}{U\theta}
\left(\frac{U}{\pi}\right)^p
\left[
K^{-p}
+
\frac{K^{1-p}}{p-1}
\right].
}
\]

This first interval backend therefore requires \(r>2\).  Lower-rank positive
forms can use a different tail bound and are not claimed by this theorem.

\subsection{Final enclosure theorem}

Let \(\boldsymbol S_K\) be the directed-rounding interval containing the finite
\(K\)-mode tilted COS sum, and set
\[
E=E_{\rm support}+E_{\rm series}.
\]
Then
\[
\boxed{
S_\theta(q)
\in
\boldsymbol S_K+[-E,E].
}
\]

If the lower endpoint of this interval is positive, directed rounding gives
\[
\boxed{
\log P(Q\ge q)
\in
K(\theta)-\theta q
+
\log\left(
\boldsymbol S_K+[-E,E]
\right).
}
\]

All finite arithmetic in this expression is enclosed by the interval backend.

\subsection{Independent exact interval control through \texorpdfstring{\(10^{-5000}\)}{1e-5000}}

For the control
\[
Q=0.02\chi^2_{80},
\]
the survival probability has the finite exact formula
\[
P(Q\ge q)
=
e^{-x}
\sum_{j=0}^{39}
\frac{x^j}{j!},
\qquad
x=\frac{q}{0.04}.
\]

This finite expression was evaluated independently by directed-rounding
interval arithmetic, without using K-COS.  The K-COS enclosure and this exact
control enclosure overlap at
\[
p=10^{-100},
\qquad
10^{-1000},
\qquad
10^{-5000}.
\]

With \(K=128\) and 90 interval decimal digits, the resulting K-COS widths are

\begin{center}
\begin{tabular}{@{}rrr@{}}
\toprule
target & width in $\log_{10}p$ & multiplicative ratio minus one\\
\midrule
$10^{-100}$ &
$8.95\times10^{-13}$ &
$2.06\times10^{-12}$\\
$10^{-1000}$ &
$6.94\times10^{-12}$ &
$1.60\times10^{-11}$\\
$10^{-5000}$ &
$3.37\times10^{-11}$ &
$7.75\times10^{-11}$\\
\bottomrule
\end{tabular}
\end{center}

Thus the \(10^{-5000}\) result is not merely a point estimate agreeing with a
reference.  For the exact supplied spectrum, it comes with a directed-rounding
logarithmic enclosure whose probability ratio satisfies
\[
\boxed{
\frac{p_+}{p_-}
\le
1+7.75\times10^{-11}
}
\]
up to the displayed rounding of that ratio.

\begin{figure}[H]
\centering
\includegraphics[width=.82\linewidth]{figure_24_interval_deep_tail_width.png}
\caption{Width of the directed-rounding K-COS probability enclosure on exact
scaled-chi-square controls.  The enclosure remains narrow even at
$10^{-5000}$.}
\label{fig:interval-tail-width}
\end{figure}

\subsection{Arithmetic width versus analytic width}

At the \(10^{-5000}\) control, the interval precision was swept from 40 to 120
decimal digits.  The total certified width remained essentially
\[
3.34247\times10^{-11}
\]
decades.

Over the same sweep, the arithmetic-only width contracted from approximately
\[
7.5\times10^{-37}
\]
to
\[
6.3\times10^{-117}
\]
decades.

\begin{figure}[H]
\centering
\includegraphics[width=.82\linewidth]{figure_25_interval_arithmetic_vs_analytic.png}
\caption{Arithmetic versus analytic uncertainty at a $10^{-5000}$ tail.  The
directed-rounding arithmetic enclosure becomes negligible; the final width is
controlled by the explicit analytic support and series remainders.}
\label{fig:interval-arithmetic}
\end{figure}

This separation is important.  Increasing numerical precision cannot improve
an analytic truncation theorem.  Conversely, once interval arithmetic is far
below the analytic error budget, numerical rounding is no longer the limiting
uncertainty.

\subsection{Guarantee boundary}

The v1.4 statement is intentionally narrow:

\begin{quote}
For an exact supplied positive decimal spectrum of rank greater than two, the
backend returns a directed-rounding logarithmic probability enclosure enlarged
by proved analytic support/alias and omitted-mode bounds.
\end{quote}

The result is conditional on correctness of the interval-arithmetic library
implementation.  The current implementation uses \texttt{mpmath.iv}, whose
interval layer applies downward and upward endpoint rounding.

A completely end-to-end guarantee beginning from an arbitrary floating-point
matrix requires one more layer: either a rigorously enclosed matrix-to-spectrum
computation, or an exact/residual completion step whose linear algebra is also
performed with verified interval or ball arithmetic.

The current packetized spectral representation has likewise not yet been
ported to interval arithmetic.  The first rigorous backend deliberately uses
the exact supplied spectrum so that representation error cannot be hidden
inside the numerical enclosure.


\section{Accuracy-driven certification}\label{sec:final-5}

The public numerical contract should specify the desired accuracy, not the
internal transform resolution.

For an exact positive spectrum and threshold \(q\), the solver accepts a
requested relative multiplicative accuracy
\[
\varepsilon_{\rm req}>0
\]
and returns a certified interval
\[
[p_-,p_+]
\]
only when
\[
\boxed{
\frac{p_+}{p_-}-1
\le
\varepsilon_{\rm req}.
}
\]

In logarithmic coordinates, if
\[
\log_{10}p\in[d_-,d_+],
\]
the acceptance rule is
\[
\boxed{
10^{d_+-d_-}-1
\le
\varepsilon_{\rm req}.
}
\]

No value of \(K\), support width, or arithmetic precision appears in the
external contract.

\begin{theorem}[Accuracy contract for the positive exact-spectrum backend]\label{claim:final-5-1}
Assume an exact supplied positive decimal spectrum of rank greater than two,
a threshold \(q\ge0\) with valid interval-design parameters, and requested multiplicative tolerance
\(\varepsilon_{\rm req}>0\).  Suppose the directed-rounding interval backend
returns an enclosure
\[
\log_{10}p\in[d_-,d_+]
\]
that includes the proved support/alias and omitted-mode remainders.  If the
accuracy-driven controller returns \texttt{certified=True}, then
\[
\boxed{
P(Q\ge q)\in[10^{d_-},10^{d_+}]
}
\]
and
\[
\boxed{
\frac{10^{d_+}}{10^{d_-}}-1
\le
\varepsilon_{\rm req}.
}
\]
If the resource caps are reached before this inequality is proved, the
controller returns \texttt{certified=False} rather than promoting the interval
midpoint to a certified answer.
\end{theorem}

\begin{proof}
The directed-rounding backend encloses the finite transform and exponential
prefactor and then enlarges that enclosure by the proved support/alias and
infinite-mode error bounds.  Therefore the true tail probability lies inside
the returned interval.  The controller encloses
\(\operatorname{expm1}(\log(10)(d_+-d_-))\) by directed rounding and
sets \texttt{certified=True} only when the upper endpoint of that ratio
enclosure is no larger than the lower endpoint of the exact requested
tolerance enclosure.  High precision alone is not an enclosure proof.  The fail-closed branch follows directly from
the controller definition.
\end{proof}

\subsection{Adaptive refinement}

The current controller adjusts three internal resources:

\begin{enumerate}
\item the COS mode count \(K\);
\item the tilted support scale \(L\), which changes \(U\) and its associated
Chernoff parameter;
\item the interval arithmetic precision.
\end{enumerate}

At each stage the full v1.4 directed-rounding interval is recomputed.  If the
rounding component is material relative to the requested budget, arithmetic
precision is increased.  Otherwise the dominant analytic contribution
determines the next refinement: mode error increases \(K\), while
support/alias error increases \(L\).

The solver terminates only after the multiplicative inequality above has been
proved.

\subsection{Fail-closed behavior}

If the configured resource caps are exhausted before the requested accuracy
is established, the solver does not silently return its midpoint as though it
were certified.

Instead it returns
\[
\boxed{\texttt{certified = false}}
\]
together with the unresolved interval and a failure reason.

This behavior is part of the numerical definition of the production solver.

\subsection{Accuracy ladder through \texorpdfstring{\(10^{-5000}\)}{1e-5000}}

The exact control
\[
Q=0.02\chi^2_{80}
\]
was evaluated at tail depths
\[
10^{-100},
\qquad
10^{-1000},
\qquad
10^{-5000}.
\]

For each depth the controller was asked to prove relative multiplicative
accuracies
\[
10^{-8},
\quad
10^{-10},
\quad
10^{-12},
\quad
10^{-14}.
\]
All twelve contracts were satisfied.

At the \(10^{-5000}\) point, the controller produced:

\begin{center}
\begin{tabular}{@{}rrrr@{}}
\toprule
requested accuracy & achieved ratio minus one & \(K\) & \(L\)\\
\midrule
\(10^{-8}\)  & \(3.19\times10^{-10}\) & 64  & 10\\
\(10^{-10}\) & \(7.70\times10^{-11}\) & 128 & 10\\
\(10^{-12}\) & \(3.48\times10^{-15}\) & 128 & 12\\
\(10^{-14}\) & \(3.48\times10^{-15}\) & 128 & 12\\
\bottomrule
\end{tabular}
\end{center}

The gap between requested and achieved accuracy is expected: refinement is
discrete, so one additional support or mode refinement can improve the
certificate by several orders of magnitude.

A further stress test requested
\[
10^{-16}
\quad\text{and}\quad
10^{-18}
\]
relative accuracy at the same \(10^{-5000}\) tail.  The solver obtained
approximately
\[
3.05\times10^{-18}
\]
and
\[
9.20\times10^{-20}
\]
respectively.

\begin{figure}[H]
\centering
\includegraphics[width=.82\linewidth]{figure_26_accuracy_driven_refinement.png}
\caption{Accuracy-driven refinement.  Each point lies at or below the
requested-equals-achieved line because an interval result is labeled certified only
after the requested multiplicative probability accuracy has been proved.}
\label{fig:accuracy-driven}
\end{figure}

\subsection{Acceptance arithmetic and serialization are directed}

The acceptance condition must be evaluated as carefully as the probability
interval.

An early implementation converted a tiny log-width \(w\) using
\[
10^w-1
\]
under the default multiprecision context.  At approximately \(10^{-15}\)
relative scales this could round the ratio to zero.

The enclosure itself remained valid, but that conversion was not acceptable
for a certification decision.  The final-edition implementation encloses
\[
\boxed{
\operatorname{expm1}(\log(10)\,w)
}
\]
with directed interval arithmetic before deciding whether the requested
accuracy has been met.  Acceptance compares the upper endpoint of that ratio
with the lower endpoint of the requested exact tolerance, so a rounded point
estimate cannot promote an unresolved boundary case to success.

Authoritative output strings must preserve the same outer enclosure.  The
final edition converts each exact binary interval endpoint to a decimal
string with \texttt{Decimal} and explicit \texttt{ROUND\_FLOOR} for lower
endpoints or \texttt{ROUND\_CEILING} for upper endpoints.  Ordinary interval
display formatting is not used as an authoritative serialization method:
nearest-rounded decimal text can move an endpoint inward even when the
internal interval arithmetic was directed.  Matrix-to-tail propagation keeps
these endpoint strings and applies the same directed ratio test.

This regression is retained because the correctness requirement applies not
only to the tail integral but also to the logic that labels the answer
certified.

\subsection{Input semantics at \texorpdfstring{\(p\approx10^{-16}\)}{p about 1e-16}}

At this accuracy scale, the representation of the threshold \(q\) itself is
measurable.  In the \(0.02\chi^2_{80}\) control targeted at \(10^{-16}\),
rounding the high-precision threshold to binary float64 shifts the exact tail
by approximately
\[
3.83\times10^{-15}
\]
decades.  The hardened float64 solver evaluated at that same rounded input
then differs from the exact tail of the rounded input by approximately
\[
5.17\times10^{-15}
\]
decades.

Accordingly, claims below the float64 input-quantization scale use exact-decimal
or interval inputs.  The rigorous accuracy contract therefore certifies the
mathematical inputs supplied to it, rather than silently identifying a rounded
binary threshold with an ideal decimal threshold.


\section{Verified matrix-to-spectrum bridge}\label{sec:final-6}

The strongest tail result above begins from a positive spectrum.  This section
moves the rigor boundary one level upstream by certifying the spectrum produced
from a floating-point factor matrix.

\subsection{Binary64 input semantics}

Let
\[
A\in\mathbb R^{n\times m}
\]
be supplied as an IEEE binary64 array.  Each stored entry is treated as its
exact dyadic rational value.  The target operator is therefore the exact Gram
matrix
\[
H=AA^T
\]
defined by those binary64 values, not an unspecified nearby real matrix.

A high-precision symmetric eigensolver is used only to generate candidate
quantities
\[
Q\in\mathbb R^{n\times n},
\qquad
D=\operatorname{diag}(d_1,\ldots,d_n),
\qquad
d_1\ge\cdots\ge d_n\ge0.
\]
No correctness claim depends on trusting this eigensolver.

\subsection{Post-certified spectral enclosure theorem}

\begin{theorem}[Candidate eigensystem post-certification]\label{claim:final-6-1}
Let
\[
B=QDQ^T.
\]
Suppose directed-rounding computation proves
\[
\|Q^TQ-I\|_2\le\eta<1
\]
and
\[
\|H-B\|_2\le\rho.
\]
Then the ordered eigenvalues of the positive semidefinite matrix \(H\) satisfy
\[
\boxed{
\max\{0,(1-\eta)d_i-\rho\}
\le
\lambda_i(H)
\le
(1+\eta)d_i+\rho
}
\]
for every \(i\).
\end{theorem}

\begin{proof}
Set
\[
S=Q^TQ.
\]
The first hypothesis gives the Loewner inequalities
\[
(1-\eta)I\preceq S\preceq(1+\eta)I.
\]
The eigenvalues of
\[
B=(QD^{1/2})(QD^{1/2})^T
\]
are the same as those of
\[
D^{1/2}SD^{1/2}.
\]
Hence
\[
(1-\eta)D
\preceq
D^{1/2}SD^{1/2}
\preceq
(1+\eta)D,
\]
and eigenvalue monotonicity yields
\[
(1-\eta)d_i
\le
\lambda_i(B)
\le
(1+\eta)d_i.
\]
Weyl's perturbation inequality and
\(\|H-B\|_2\le\rho\) then give
\[
|\lambda_i(H)-\lambda_i(B)|\le\rho.
\]
Combining the two bounds and using \(H\succeq0\) proves the claim.
\end{proof}

The implementation does not attempt to evaluate the spectral norms above
exactly.  Instead it encloses every entry of
\[
Q^TQ-I
\quad\text{and}\quad
H-QDQ^T
\]
with directed-rounding interval arithmetic and uses
\[
\|M\|_2\le\|M\|_F.
\]
Thus the reported \(\eta\) and \(\rho\) are conservative certified upper
bounds.

\subsection{Monotone propagation into the tail}

Let the spectral theorem produce
\[
\ell_i\le\lambda_i(H)\le u_i.
\]
Using the same independent standard normals in all three sums gives
\[
\sum_i\ell_iZ_i^2
\le
\sum_i\lambda_iZ_i^2
\le
\sum_i u_iZ_i^2
\qquad\text{almost surely}.
\]
Therefore
\[
\boxed{
P\!\left(\sum_i\ell_iZ_i^2\ge q\right)
\le
P(Q\ge q)
\le
P\!\left(\sum_i u_iZ_i^2\ge q\right).
}
\]
The lower endpoint is evaluated by the rigorous positive-tail backend on the
lower spectrum, and the upper endpoint by the same backend on the upper
spectrum.  The current controller requires more than two strictly positive
coefficients in each endpoint spectrum and a threshold above both endpoint
means.  An endpoint spectrum outside these implemented conditions is rejected
or reported as unsupported; monotone coupling itself remains valid.  The final matrix-defined probability interval is accepted only when
its multiplicative ratio satisfies the requested accuracy.

\subsection{Random dense spectral-containment audit}

Three unrelated dense binary64 factors were tested against independently
recomputed high-precision eigenvalues of their exact dyadic Gram matrices:

\begin{center}
\begin{tabular}{@{}rrrr@{}}
\toprule
factor size & all eigenvalues contained & maximum relative width & time (s)\\
\midrule
$12\times18$ & yes & $3.36\times10^{-58}$ & 0.09\\
$20\times30$ & yes & $1.44\times10^{-57}$ & 0.35\\
$28\times36$ & yes & $1.20\times10^{-56}$ & 0.80\\
\bottomrule
\end{tabular}
\end{center}

No reference eigenvalue fell outside its certified interval.

\subsection{End-to-end matrix-defined tail near \texorpdfstring{\(10^{-16}\)}{1e-16}}

A dense \(64\times64\) exact-binary factor was constructed from a normalized
Hadamard basis with four singular-value groups.  The stored binary64 entries
are exact dyadic rationals.  The spectral postcheck returned
\[
\eta\le1.65\times10^{-49},
\qquad
\rho\le9.18\times10^{-50}.
\]
For example, the leading coefficient was enclosed by
\[
0.999\ldots97435
<
\lambda_1
<
1.000\ldots02565,
\]
with a width of order \(10^{-49}\).

At
\[
q=131.33201485100517,
\]
the tail lies near \(10^{-16}\).  Starting from the factor matrix itself:

\begin{center}
\begin{tabular}{@{}rrr@{}}
\toprule
requested relative accuracy & achieved ratio minus one & certified\\
\midrule
$10^{-10}$ & $2.60\times10^{-12}$ & yes\\
$10^{-12}$ & $4.09\times10^{-15}$ & yes\\
\bottomrule
\end{tabular}
\end{center}

For the \(10^{-12}\) request, the final enclosure was approximately
\[
\log_{10}p
\in
[-15.99999999999996625,
 -15.99999999999996447].
\]

This closes the basic numerical chain
\[
\boxed{
\text{binary64 factor}
\longrightarrow
\text{certified spectrum}
\longrightarrow
\text{certified deep-tail interval}
}
\]
at moderate matrix dimension.

\subsection{Scalability boundary}

The theorem is general, but the present verifier forms a full candidate
basis \(Q\) and postchecks dense matrices.  It therefore closes the correctness
boundary before it closes the large-scale memory boundary.  For very large
streamed factors, the next step is to replace the full-basis postcheck with a
verified range-core/residual certificate and interval exact-completion step,
so that the matrix-to-spectrum guarantee inherits the same memory-bounded
architecture developed elsewhere in this manuscript.


\section{Weighted-feature specialization}\label{sec:final-7}

A prominent application constructs the general factor from weighted,
covariate-residualized features.  A rare-variant kernel statistic is one concrete instance, but the same
formulas apply to any feature matrix with nonnegative column weights.

Let
\[
G\in\mathbb R^{n\times m},
\qquad
W=\operatorname{diag}(w_1,\ldots,w_m),
\quad w_j\ge 0,
\]
and let \(Q_X\in\mathbb R^{n\times p}\) have orthonormal columns spanning
the null-model covariates.  Define
\[
M_X=I-Q_XQ_X^\top,
\qquad
A=M_XGW^{1/2},
\qquad
K=AA^\top.
\]

The central implementation difficulty is not the algebraic size of the
model but the mismatch between the total factor size \(n\times m\) and the
working memory of an accelerator.  The key point of this manuscript is that
several decompositions act on different mathematical coordinates and therefore
compose:
\[
\boxed{
\text{storage blocks}
\circ
\text{structural phases}
\circ
\text{Krylov/range compression}
\circ
\text{spectral tail packets}.
}
\]

Storage blocking controls residence.  Structural phases encode LD windows,
annotations, frequency bands, or other coordinate groups.  Krylov methods
compress the operator to the rank actually requested.  Spectral packets
compress the final generalized-\(\chi^2\) transform.


\begin{figure}[H]
\centering
\includegraphics[width=.72\linewidth]{figure_01_residualization_geometry.png}
\caption{Residualization is the geometric starting point. Every sample-space vector splits into a component in the covariate span and an orthogonal residual component. The factor $A=M_XGW^{1/2}$ therefore acts only on variation remaining after the null-model covariates have been removed.}
\label{fig:residualization-geometry}
\end{figure}

\section{Variant-block Gram decomposition}\label{sec:final-8}

Partition the variant indices into disjoint blocks
\[
\{1,\ldots,m\}
=
J_1\sqcup\cdots\sqcup J_B.
\]
Let
\[
G_b=G_{[:,J_b]},
\qquad
W_b=W_{J_b,J_b},
\qquad
A_b=M_XG_bW_b^{1/2}.
\]
Then
\[
A=[A_1\ \cdots\ A_B].
\]

\begin{theorem}[Additive block Gram identity]\label{claim:final-8-1}
\[
\boxed{
K
=
AA^\top
=
\sum_{b=1}^B A_bA_b^\top
=
M_X
\left(
\sum_{b=1}^B G_bW_bG_b^\top
\right)
M_X.
}
\]
\end{theorem}

\begin{proof}
Horizontal concatenation gives
\[
[A_1\ \cdots\ A_B]
\begin{bmatrix}
A_1^\top\\
\vdots\\
A_B^\top
\end{bmatrix}
=
\sum_bA_bA_b^\top.
\]
The final form follows from
\(A_b=M_XG_bW_b^{1/2}\) and
\(M_X^\top=M_X=M_X^2\).
\end{proof}

\begin{remark}\label{claim:final-8-2}
The theorem is exact for every partition.  Block boundaries are execution
parameters, not statistical parameters.
\end{remark}


\begin{figure}[H]
\centering
\includegraphics[width=.92\linewidth]{figure_02_exact_block_kernel_assembly.png}
\caption{Exact block-kernel assembly. The sample interaction kernel carried by each resident variant block is $K_b=A_bA_b^\top$, and the full kernel is their exact sum. Streaming changes the order of evaluation, not the operator.}
\label{fig:block-kernel-assembly}
\end{figure}

\section{Exact streamed kernel actions}\label{sec:final-9}

\begin{theorem}[One-pass sample-space streaming action]\label{claim:final-9-1}
For every block vector \(V\in\mathbb R^{n\times s}\),
\[
\boxed{
KV
=
M_X
\sum_{b=1}^B
G_bW_bG_b^\top
M_XV.
}
\]
\end{theorem}

\begin{proof}
Apply the additive block Gram identity to \(KV\), then associate each block
product as
\[
G_bW_b(G_b^\top M_XV).
\]
\end{proof}

\begin{algorithm}[Streamed \(KV\)]\label{claim:final-9-2}
\begin{enumerate}[label=\arabic*.]
\item Compute \(U=M_XV\).
\item Set \(Y=0\).
\item For each resident block \(b\), compute
\[
T_b=G_b^\top U,
\qquad
Y\leftarrow Y+G_b(W_bT_b).
\]
\item Return \(M_XY\).
\end{enumerate}
\end{algorithm}

\begin{corollary}[Working memory independent of total variant count]\label{claim:final-9-3}
If at most \(c\) variants are resident and the block-vector width is \(s\),
then the algebraic working set is
\[
\boxed{
O(nc+ns+cs+np),
}
\]
independent of total variant count \(m\).
\end{corollary}

\section{Streamed factor actions and Golub--Kahan}\label{sec:final-10}

Let \(V\in\mathbb R^{m\times s}\) be partitioned conformally:
\[
V=
\begin{bmatrix}
V_1\\
\vdots\\
V_B
\end{bmatrix}.
\]

\begin{theorem}[Streamed \(AV\)]\label{claim:final-10-1}
\[
\boxed{
AV
=
M_X
\sum_{b=1}^B
G_bW_b^{1/2}V_b.
}
\]
\end{theorem}

\begin{theorem}[Streamed \(A^\top U\)]\label{claim:final-10-2}
For \(U\in\mathbb R^{n\times s}\),
\[
\boxed{
(A^\top U)_b
=
W_b^{1/2}G_b^\top M_XU.
}
\]
\end{theorem}

\begin{proof}
Both identities follow by substituting
\(A=M_XGW^{1/2}\) and partitioning the variant coordinate.
\end{proof}

\begin{corollary}[Streaming Golub--Kahan equivalence]\label{claim:final-10-3}
Any exact-arithmetic algorithm whose access to \(A\) occurs only through
\(AV\) and \(A^\top U\) produces the same iterates when those products are
implemented by the preceding streamed formulas.
\end{corollary}

This applies directly to block Golub--Kahan bidiagonalization, matrix-free
singular-value iteration, range finding, and restarted variants.  The
factor need not be resident.

\section{Two-pass dual action}\label{sec:final-11}

The variant-space normal operator is
\[
H=A^\top A.
\]

\begin{theorem}[Two-pass streamed dual action]\label{claim:final-11-1}
For every \(V\in\mathbb R^{m\times s}\),
\[
\boxed{
HV=A^\top(AV).
}
\]
The product is evaluated exactly by a streamed pass forming
\(Y=AV\in\mathbb R^{n\times s}\), followed by a second streamed pass
emitting the blocks
\[
(HV)_b=W_b^{1/2}G_b^\top M_XY.
\]
\end{theorem}

This form is particularly natural for variant-space polyphase Krylov
iterations because the right-coordinate blocks can be produced and consumed
sequentially.

\section{Two-dimensional sample/variant tiling}\label{sec:final-12}

A variant block may itself exceed device memory.  Let
\[
H_b=G_bW_b^{1/2}
\]
and partition samples into row tiles \(I_1,\ldots,I_R\).

\begin{theorem}[Two-pass residualized tile construction]\label{claim:final-12-1}
Define
\[
R_b=Q_X^\top H_b.
\]
Then
\[
\boxed{
R_b
=
\sum_r
Q_{X,I_r}^\top H_{b,I_r},
}
\]
and after \(R_b\) is known,
\[
\boxed{
A_{b,I_r}
=
H_{b,I_r}-Q_{X,I_r}R_b.
}
\]
Consequently,
\[
\boxed{
K_{I_r,I_s}
=
\sum_b
A_{b,I_r}A_{b,I_s}^\top.
}
\]
\end{theorem}

\begin{proof}
The first identity is the row-partitioned matrix product
\(Q_X^\top H_b\).  The second is the definition
\(M_XH_b=H_b-Q_XQ_X^\top H_b\).  The last identity is the corresponding
sample-tile expansion of \(AA^\top\).
\end{proof}

With sample tile height \(t\) and resident variant width \(c\), the dominant
working storage, including a resident covariate tile, is
\[
O(tc+pc+tp+t^2).
\]
This count excludes the backing input store, completed output blocks,
allocator overhead, and library workspaces.

\section{Partition refinement and permutation invariance}\label{sec:final-13}

\begin{theorem}[Partition-refinement invariance]\label{claim:final-13-1}
Let \(\mathcal J\) be any partition of the variants and let
\(\mathcal J'\) refine \(\mathcal J\).  Then
\[
\boxed{
\sum_{J\in\mathcal J}A_JA_J^\top
=
\sum_{J'\in\mathcal J'}A_{J'}A_{J'}^\top
=
AA^\top.
}
\]
The analogous streamed operator action is also invariant.
\end{theorem}

\begin{proof}
Each coarse block is a disjoint union of refined blocks.  Apply the additive
Gram theorem within each coarse block and sum.
\end{proof}

\begin{corollary}[Adaptive memory splitting]\label{claim:final-13-2}
A runtime may recursively subdivide any block that exceeds available device
memory without changing the exact operator.
\end{corollary}

\begin{proposition}[Variant permutation invariance]\label{claim:final-13-3}
If the columns of \(G\) and the diagonal entries of \(W\) are permuted by the
same permutation, \(K\) and every streamed sample-space action are unchanged.
\end{proposition}


\begin{figure}[H]
\centering
\includegraphics[width=.84\linewidth]{figure_08_partition_invariance.png}
\caption{Partition-refinement invariance in floating-point tests. Coarse, fine, irregular, singleton-heavy, and permuted blockings reproduce the explicitly formed Gram matrix and its action to machine precision. The block schedule is therefore an execution choice rather than a change in the statistical model.}
\label{fig:partition-invariance}
\end{figure}

\section{Polyphase structure}\label{sec:final-14}

Let \(P_h\) project onto variant phase \(h\), and again set
\[
H=A^\top A.
\]
Define the exact phase graph by
\[
g\to h
\quad\Longleftrightarrow\quad
P_hHP_g\ne0.
\]

\begin{theorem}[Polyphase light cone]\label{claim:final-14-1}
If the graph distance from phase \(g\) to phase \(h\) exceeds \(k\), then
\[
\boxed{
P_hH^kP_g=0.
}
\]
\end{theorem}

\begin{proof}
Insert the projector resolution
\(\sum_jP_j=I\) between successive factors of \(H\).  Every nonzero term in
\(P_hH^kP_g\) corresponds to a phase path of length \(k\).  If no such path
exists, every term vanishes.
\end{proof}


\begin{figure}[H]
\centering
\includegraphics[width=.82\linewidth]{figure_04_polyphase_lightcone.png}
\caption{The exact polyphase light cone. The support of the Krylov iterate expands by at most one graph edge per operator application. Phases outside graph distance $k$ remain identically inactive at depth $k$ and need not be loaded.}
\label{fig:polyphase-lightcone}
\end{figure}

\begin{theorem}[Streaming-polyphase composition]\label{claim:final-14-2}
At Krylov depth \(k\), an exact matrix-free iteration may omit every phase
outside the exact light cone.  Within each active phase, arbitrary storage
refinement may be used.
\end{theorem}

\begin{proof}
The light-cone theorem removes inactive phase contributions exactly.
Partition-refinement invariance shows that the remaining active phase
operators may be streamed through any exact subpartition.
\end{proof}

This is the key composition law:
\[
\boxed{
\text{mathematical locality}
\longrightarrow
\text{active phases}
\longrightarrow
\text{active storage blocks}.
}
\]

\section{Certified approximate light cones}
\label{sec:approx-lightcone}

The exact light-cone theorem requires literal zero phase couplings.  Real LD,
annotation, and feature graphs are often only approximately sparse.  This
section replaces exact support by a computable phasewise envelope.

Let the phase projectors be $P_1,\ldots,P_q$, and let $X$ be a vector or block
vector.  Define its phase-norm profile
\[
\rho(X)_g=\lVert P_gX\rVert_F,
\]
and define the nonnegative phase-transfer matrix
\[
\mathcal B_{hg}=\lVert P_hHP_g\rVert_2.
\]
The matrix $\mathcal B$ records the strength of every one-step phase transfer
without assuming that weak couplings vanish.

\begin{theorem}[Weighted phase-transfer envelope]
\label{thm:weighted-transfer}
For every block vector $X$,
\[
\boxed{
\rho(HX)\leq \mathcal B\rho(X)
}
\]
componentwise.  Consequently, for every integer $k\geq0$,
\[
\boxed{
\rho(H^kX)\leq \mathcal B^k\rho(X).
}
\]
\end{theorem}

\begin{proof}
Using $\sum_gP_g=I$ and submultiplicativity,
\[
\begin{aligned}
\lVert P_hHX\rVert_F
&=\left\lVert\sum_gP_hHP_gX\right\rVert_F\\
&\leq\sum_g\lVert P_hHP_g\rVert_2\lVert P_gX\rVert_F\\
&=(\mathcal B\rho(X))_h.
\end{aligned}
\]
Iteration gives the power bound.
\end{proof}

Now choose a symmetric set of retained phase blocks and define
\[
H_0=\sum_{(h,g)\in\mathcal E_0}P_hHP_g.
\]
Let
\[
(\mathcal B_0)_{hg}=
\begin{cases}
\mathcal B_{hg},&(h,g)\in\mathcal E_0,\\
0,&\text{otherwise},
\end{cases}
\qquad
\mathcal R=\mathcal B-\mathcal B_0.
\]
Symmetric pair retention preserves self-adjointness of $H_0$, although the
following error theorem does not require positive semidefiniteness of $H_0$.

\begin{theorem}[Thresholded-power certificate]
\label{thm:thresholded-power}
For every $k\geq1$,
\[
\boxed{
\rho\bigl((H^k-H_0^k)X\bigr)
\leq
(\mathcal B^k-\mathcal B_0^k)\rho(X).
}
\]
In particular,
\[
\boxed{
\lVert(H^k-H_0^k)X\rVert_F
\leq
\left\lVert(\mathcal B^k-\mathcal B_0^k)\rho(X)\right\rVert_2.
}
\]
\end{theorem}

\begin{proof}
Write $E=H-H_0$.  The telescoping identity gives
\[
H^k-H_0^k
=
\sum_{j=0}^{k-1}H^{k-1-j}EH_0^j.
\]
The block-norm matrix of $E$ is $\mathcal R$.  Applying
Theorem~\ref{thm:weighted-transfer} successively to each term gives
\[
\rho\bigl((H^k-H_0^k)X\bigr)
\leq
\sum_{j=0}^{k-1}
\mathcal B^{k-1-j}\mathcal R\mathcal B_0^j\rho(X).
\]
The matrix telescoping identity
\[
\mathcal B^k-\mathcal B_0^k
=
\sum_{j=0}^{k-1}
\mathcal B^{k-1-j}(\mathcal B-\mathcal B_0)\mathcal B_0^j
\]
yields the phasewise formula.  Orthogonality of the phase projectors gives
\[
\lVert Y\rVert_F^2=\sum_h\lVert P_hY\rVert_F^2,
\]
which proves the global bound.
\end{proof}

\begin{corollary}[Certified approximate light cone]
\label{cor:approx-lightcone}
Assume $X$ is supported on a set of start phases $S$.  If phase $h$ is not
reachable from $S$ by a retained path of length at most $k$, then
\[
P_hH_0^kX=0
\]
and
\[
\boxed{
\lVert P_hH^kX\rVert_F
\leq
\bigl[(\mathcal B^k-\mathcal B_0^k)\rho(X)\bigr]_h.
}
\]
Thus an exact zero outside the retained light cone is replaced by a certified
nonzero leakage budget.
\end{corollary}

\begin{corollary}[Scalar threshold law]\label{claim:final-15-4}
For any induced matrix norm, let
\[
M=\max\{\lVert\mathcal B\rVert,\lVert\mathcal B_0\rVert\}.
\]
Then
\[
\boxed{
\lVert\mathcal B^k-\mathcal B_0^k\rVert
\leq
kM^{k-1}\lVert\mathcal R\rVert.
}
\]
Hence weak discarded coupling produces a first-order error bound proportional
to the discarded transfer norm and to the Krylov depth.
\end{corollary}

\begin{theorem}[Krylov-polynomial certificate]
\label{thm:polynomial-certificate}
Let
\[
p(t)=\sum_{j=0}^d a_jt^j.
\]
Then
\[
\boxed{
\rho\bigl((p(H)-p(H_0))X\bigr)
\leq
\sum_{j=1}^d
|a_j|(\mathcal B^j-\mathcal B_0^j)\rho(X).
}
\]
\end{theorem}

\begin{proof}
Apply Theorem~\ref{thm:thresholded-power} to every nonconstant power and use
the triangle inequality.  The constant terms cancel.
\end{proof}

The polynomial theorem covers arbitrary vectors in a finite Krylov space,
including polynomial filters and restarted combinations, rather than only the
monomial $H^kX$.

\begin{corollary}[Certified phase pruning]\label{claim:final-15-6}
If a thresholded phase graph is chosen so that
\[
\left\lVert
(\mathcal B^k-\mathcal B_0^k)\rho(X)
\right\rVert_2
\leq\tau,
\]
then replacing $H^kX$ by $H_0^kX$ incurs Frobenius error at most $\tau$.
The same statement holds for a Krylov polynomial with the certificate in
Theorem~\ref{thm:polynomial-certificate}.
\end{corollary}

This turns graph thresholding into an error-budget decision.  A phase or
interface is not discarded because it merely appears small; it is discarded
only when the weighted sum of all paths using discarded couplings remains
inside the requested tolerance.

\begin{figure}[H]
\centering
\includegraphics[width=.84\linewidth]{figure_09_approximate_lightcone_certificate.png}
\caption{Weak long-range phase couplings replace the exact zero outside the retained light cone by a nonzero but certified envelope.  In this PSD test, the phasewise transfer-matrix certificate bounds the observed leakage at every depth from one through eight.}
\label{fig:approx-lightcone-certificate}
\end{figure}

\begin{figure}[H]
\centering
\includegraphics[width=.84\linewidth]{figure_10_threshold_tradeoff.png}
\caption{Certified phase-graph sparsification at Krylov depth five.  Retaining fewer off-diagonal phase blocks reduces communication and operator work while the transfer-matrix bound remains above the measured truncation error.  The bound is intentionally conservative because it ignores cancellation and singular-vector alignment.}
\label{fig:threshold-tradeoff}
\end{figure}

\subsection{Numerical audit}

The validation operator has twelve phases of width four.  A block-bidiagonal
factor generates a positive-definite local operator $H_0$, and weak symmetric
long-range blocks are then added to obtain $H$.  Both normalized operators
remain positive definite.  A three-column start block is supported in phase
zero.

For depths one through eight:
\begin{itemize}
\item every observed phase norm lay below $\mathcal B^k\rho(X)$;
\item every observed thresholding error lay below
$(\mathcal B^k-\mathcal B_0^k)\rho(X)$;
\item the largest total observed-error-to-certificate ratio was approximately
$0.669$.
\end{itemize}

A separate hard-threshold sweep retained between $22$ and $132$ directed
off-diagonal phase blocks.  Every measured depth-five error remained below its
certificate; the largest observed-error-to-certificate ratio was approximately
$0.376$.  Finally, a degree-four signed polynomial in $H$ satisfied the
phasewise polynomial certificate in all twelve phases, with maximum ratio
approximately $0.068$.

These computations validate the implementation and illustrate sharpness.  The
universal result is supplied by the preceding proofs, not by the finite test.


\section{Certified budget-optimal phase scheduling}
\label{sec:budget-scheduler}

The preceding section answers a validation question: given a retained phase
operator $H_0$, how large can the Krylov error be?  This section turns that
certificate into a selection rule.  The goal is to remove as much phase
communication as possible while remaining inside a prescribed error budget.

Let $\mathcal E$ denote the set of undirected off-diagonal phase pairs.  Because
$H$ is self-adjoint, $\mathcal B$ is symmetric.  For a phase pair
$e=\{h,g\}$ define the nonnegative matrix $\Delta_e$ by
\[
(\Delta_e)_{hg}=\mathcal B_{hg},\qquad
(\Delta_e)_{gh}=\mathcal B_{gh},
\]
with every other entry zero.  If $\mathcal O\subseteq\mathcal E$ is the set
of omitted pairs, define
\[
\mathcal B_{\mathcal O}
=\mathcal B-\sum_{e\in\mathcal O}\Delta_e.
\]
The diagonal is always retained.

For depth $k\ge1$ and starting phase profile $\rho_0=\rho(X)$, define the
\emph{edge risk vector}
\[
\boxed{
q_e^{(k)}
=
\sum_{j=0}^{k-1}
\mathcal B^{k-1-j}\Delta_e\mathcal B^j\rho_0
}
\]
and its scalar risk
\[
\boxed{
\beta_e^{(k)}=\lVert q_e^{(k)}\rVert_2.
}
\]
Unlike the raw block norm $\mathcal B_{hg}$, this quantity depends on the
start profile and on Krylov depth.

\begin{theorem}[Edge-additive leakage majorant]
\label{thm:edge-additive-majorant}
For every omitted set $\mathcal O\subseteq\mathcal E$,
\[
\boxed{
0\le
(\mathcal B^k-\mathcal B_{\mathcal O}^k)\rho_0
\le
\sum_{e\in\mathcal O}q_e^{(k)}
}
\]
componentwise.  Consequently,
\[
\boxed{
\lVert(H^k-H_{\mathcal O}^k)X\rVert_F
\le
\left\lVert
(\mathcal B^k-\mathcal B_{\mathcal O}^k)\rho_0
\right\rVert_2
\le
\sum_{e\in\mathcal O}\beta_e^{(k)}.
}
\]
\end{theorem}

\begin{proof}
The first inequality in the second display is
Theorem~\ref{thm:thresholded-power}.  For the transfer matrices, the exact
noncommutative telescoping identity is
\[
\mathcal B^k-\mathcal B_{\mathcal O}^k
=
\sum_{j=0}^{k-1}
\mathcal B^{k-1-j}
(\mathcal B-\mathcal B_{\mathcal O})
\mathcal B_{\mathcal O}^j.
\]
All matrices and $\rho_0$ are nonnegative, and
$0\le\mathcal B_{\mathcal O}\le\mathcal B$ entrywise.  Hence
\[
\mathcal B_{\mathcal O}^j\rho_0
\le
\mathcal B^j\rho_0.
\]
Substituting
\[
\mathcal B-\mathcal B_{\mathcal O}
=
\sum_{e\in\mathcal O}\Delta_e
\]
and distributing the sum gives the componentwise majorant.  The final scalar
bound follows from the Euclidean triangle inequality.
\end{proof}

\begin{corollary}[Polynomial edge-risk certificate]\label{claim:final-16-2}
Let
\[
p(t)=\sum_{d=0}^{D}a_dt^d.
\]
Define
\[
q_e^{(p)}
=
\sum_{d=1}^{D}|a_d|q_e^{(d)},
\qquad
\beta_e^{(p)}=\lVert q_e^{(p)}\rVert_2.
\]
Then
\[
\boxed{
\lVert(p(H)-p(H_{\mathcal O}))X\rVert_F
\le
\sum_{e\in\mathcal O}\beta_e^{(p)}.
}
\]
Thus the same scheduling machinery applies to finite Krylov combinations and
polynomial filters.
\end{corollary}

\begin{corollary}[Query-invisible phase pair]\label{claim:final-16-3}
If
\[
\Delta_e\mathcal B^j\rho_0=0
\qquad
\text{for every }0\le j<k,
\]
then
\[
\beta_e^{(k)}=0.
\]
Therefore a phase pair can be large in operator norm and still have zero
risk for a particular finite-depth query.  A magnitude-only threshold cannot
represent this distinction.
\end{corollary}

\subsection{Equal-cost optimality}

Suppose each retained phase pair has the same communication or scheduling
cost.  Under the additive certificate, a set of omitted pairs is feasible if
\[
\sum_{e\in\mathcal O}\beta_e^{(k)}\le\tau.
\]
Sort the risks as
\[
\beta_{(1)}\le\beta_{(2)}\le\cdots\le\beta_{(N)}.
\]

\begin{theorem}[Cardinality-optimal certified pruning]
\label{thm:unit-cost-optimal}
Let $r$ be the largest integer satisfying
\[
\sum_{i=1}^{r}\beta_{(i)}\le\tau.
\]
Then omitting the $r$ least-risk phase pairs is feasible, and no feasible
omitted set can contain more than $r$ pairs under the scalar additive
certificate.
\end{theorem}

\begin{proof}
Among all subsets of cardinality $r+1$, the smallest possible sum of risks is
\[
\sum_{i=1}^{r+1}\beta_{(i)}.
\]
By maximality of $r$, that sum exceeds $\tau$.  Hence no set of $r+1$ omitted
pairs is feasible.  The prefix of length $r$ is feasible by construction.
\end{proof}

This is a genuine optimization theorem, not a heuristic threshold rule.  The
optimization target is the additive certificate itself; the tighter exact
transfer envelope can leave additional slack.

\subsection{Exact-envelope refinement}

After the cardinality-optimal additive schedule is found, recomputing
\[
\left\lVert
(\mathcal B^k-\mathcal B_{\mathcal O}^k)\rho_0
\right\rVert_2
\]
may reveal unused budget.  Additional pairs can then be removed one at a time,
provided the exact transfer envelope remains below $\tau$.

\begin{proposition}[Certified refinement]\label{claim:final-16-5}
Any sequence of additional omissions accepted only when
\[
\left\lVert
(\mathcal B^k-\mathcal B_{\mathcal O}^k)\rho_0
\right\rVert_2\le\tau
\]
remains a valid $\tau$-certificate for the original operator error.
\end{proposition}

The refinement rule used in the reference implementation chooses, among all
currently feasible next omissions, the pair that produces the smallest
increase in the exact transfer envelope.  It is certified but is not claimed
globally optimal.

\subsection{Unequal communication costs}

Let $c_e\ge0$ denote the communication or storage cost saved by omitting pair
$e$.  Under the scalar additive certificate, the exact discrete selection
problem is
\[
\boxed{
\max_{x_e\in\{0,1\}}
\sum_e c_ex_e
\quad\text{subject to}\quad
\sum_e\beta_e^{(k)}x_e\le\tau.
}
\]
This is the $0$--$1$ knapsack problem.  Hence unequal-cost scheduling is
combinatorial in general.  The implementation supplies a safe
cost-to-risk greedy heuristic; its output remains certified because it never
exceeds the additive risk budget.  A fractional-knapsack relaxation supplies
an upper bound on the maximum possible saved cost.

\begin{figure}[H]
\centering
\includegraphics[width=.84\linewidth]{figure_11_budget_scheduler_comparison.png}
\caption{Depth-five pruning under a fixed certified error budget.  A global block-norm threshold is nearly unable to prune this example without crossing the requested tolerance, whereas query-aware edge risks remove many phase pairs.  The exact-envelope refinement safely uses slack left by the additive majorant.}
\label{fig:budget-scheduler-comparison}
\end{figure}

\begin{figure}[H]
\centering
\includegraphics[width=.82\linewidth]{figure_13_edge_risk_vs_block_norm.png}
\caption{Raw phase-block magnitude and finite-depth query risk are different quantities.  The edge score $\beta_{hg}^{(5)}$ incorporates the start profile and every depth-five transfer path through the phase graph; the block norm $\mathcal B_{hg}$ does not.}
\label{fig:edge-risk-vs-norm}
\end{figure}

\subsection{Numerical audit}

The twelve-phase operator from Section~\ref{sec:approx-lightcone} was reused at
Krylov depth five.  There are $66$ undirected off-diagonal phase pairs.
For error budgets $10^{-3}$, $2\times10^{-3}$, $5\times10^{-3}$, and
$10^{-2}$, the additive-optimal scheduler retained respectively
$61$, $56$, $45$, and $36$ pairs.  Exact-envelope refinement reduced those
counts further to
$58$, $50$, $36$, and $30$ pairs while remaining inside budget.

In the same test, the best single global threshold on the raw block norm
retained $65$ pairs at every one of those four budgets.  The difference is
structural: the global threshold is blind to start phase and Krylov depth.

A separate six-phase problem contains $15$ candidate phase pairs, so all
$2^{15}=32768$ subsets can be enumerated.  At four distinct budgets, the
sorted-risk scheduler matched the brute-force cardinality optimum under the
additive certificate in every case.  The exact-envelope greedy refinement also
matched the brute-force exact-envelope optimum in cardinality for all four
budgets in this finite audit.  The latter agreement is empirical and is not
promoted to a theorem.

For an unequal-cost eight-phase test, costs were proportional to the bytes in
symmetric dense phase blocks.  The cost-to-risk heuristic saved $5408$ bytes
out of a $9472$-byte communication proxy while using additive risk $0.011632$
under budget $0.012$.  The fractional relaxation upper bound was $5647.8$
bytes, placing the greedy schedule within about $4.25\%$ of that relaxation
bound on this instance.

\section{Asynchronous double-buffer execution model}
\label{sec:double-buffer}

The mathematical streaming identities do not require transfers and GEMMs to be
serial.  Consider $B$ storage blocks.  Let $t_b$ be the transfer time for block
$b$ and $c_b$ its device compute time.

\begin{proposition}[Ideal two-buffer pipeline]\label{claim:final-17-1}
In a model with perfect asynchronous overlap, transfer of block $b+1$ may run
while block $b$ is being processed.  The serial one-buffer time is
\[
T_{\rm serial}=\sum_{b=1}^{B}(t_b+c_b),
\]
whereas the two-buffer pipeline time is
\[
\boxed{
T_{\rm pipe}
=
t_1
+
\sum_{b=1}^{B-1}\max\{c_b,t_{b+1}\}
+
c_B.
}
\]
For uniform times $t_b=t$ and $c_b=c$,
\[
T_{\rm pipe}=t+c+(B-1)\max\{t,c\}.
\]
\end{proposition}

\begin{proof}
After the first block arrives, every interior pipeline stage finishes only when
both the current computation and the next transfer have completed.  Their
joint stage duration is therefore the maximum of those two times.  The initial
transfer and final compute cannot be hidden and form the two boundary terms.
\end{proof}

If one resident block uses $M_{\rm block}$ bytes and all other fixed working
state uses $M_{\rm fixed}$ bytes, the corresponding memory models are
\[
M_{\rm single}=M_{\rm fixed}+M_{\rm block},
\qquad
\boxed{
M_{\rm double}=M_{\rm fixed}+2M_{\rm block}.
}
\]
The proposition is a scheduling model, not a hardware timing claim.  PCIe or
NVLink contention, kernel launch overhead, synchronization, page pinning, and
library workspaces can reduce realized overlap.

\begin{figure}[H]
\centering
\includegraphics[width=.82\linewidth]{figure_12_double_buffer_model.png}
\caption{Idealized effect of asynchronous double buffering over $64$ streamed blocks.  The largest modeled speedup occurs when transfer and computation times are comparable, because neither resource is left idle for most of the pipeline.  The additional cost is one extra resident block buffer.}
\label{fig:double-buffer-model}
\end{figure}

\section{Phase-local rank}\label{sec:final-18}

Let \(R\) be a newly generated Krylov residual and
\[
R_h=P_hR.
\]

\begin{theorem}[Minimal exact phase width]\label{claim:final-18-1}
Let
\[
r_h=\operatorname{rank}(R_h).
\]
Every phase-supported basis spanning \(\operatorname{range}(R_h)\) has at
least \(r_h\) vectors, and a rank-revealing QR or compact SVD attains exactly
\(r_h\).
\end{theorem}

This justifies nonuniform Krylov block widths by phase.


\begin{figure}[H]
\centering
\includegraphics[width=.82\linewidth]{figure_05_phase_local_rank.png}
\caption{Phase-local rank as a dimension law. The archived six-phase audit records constructed rank, numerical rank and minimum exact width of one through six, with the remaining capacity wasted by a uniform width of twelve. A rank-revealing factorization keeps the new directions. This summary replaces an annotated singular-value curve whose raw plotted vector was not archived.}
\label{fig:phase-local-rank}
\end{figure}

\section{Low-rank interface messages}\label{sec:final-19}

If an inter-phase coupling has exact factorization
\[
P_hHP_g
=
U_{hg}V_{hg}^\top
\]
of rank \(r_{hg}\), then
\[
\boxed{
P_hHP_gx
=
U_{hg}(V_{hg}^\top x).
}
\]

This representation reduces storage and communication.  It does not
automatically reduce wall-clock time unless the factors are applied through
batched or fused kernels; a Python reference implementation can lose to
dense BLAS despite lower arithmetic and storage.

\section{Exact QR-core updates}\label{sec:final-20}

Suppose
\[
A=QR,
\qquad
Q^\top Q=I,
\]
and append a new block \(U\).  Set
\[
C=Q^\top U,
\qquad
F=U-QC,
\qquad
F=Q_\perp R_\perp,
\qquad Q_\perp^\top Q_\perp=I,
\qquad Q^\top Q_\perp=0.
\]
Here the columns of \(Q_\perp\) form an orthonormal basis of the nonzero
residual range; zero residual directions are omitted.

\begin{theorem}[Exact QR-core update]\label{claim:final-20-1}
\[
[A,U]
=
[Q,Q_\perp]
\begin{pmatrix}
R&C\\
0&R_\perp
\end{pmatrix}.
\]
Therefore the nonzero eigenvalues of
\([A,U][A,U]^\top\) are exactly the eigenvalues of the small core Gram.
\end{theorem}

\begin{proof}
The first block product is \(QR=A\); the second is
\(QC+Q_\perp R_\perp=QC+F=U\).  The stated orthogonality
conditions make \([Q,Q_\perp]\) an isometry on the compact core
space.  Left multiplication by this isometry preserves all nonzero singular
values of the core, and squaring them yields the Gram eigenvalues.
\end{proof}

This is useful for overlapping genomic windows and small incremental variant
updates.


\begin{figure}[H]
\centering
\includegraphics[width=.84\linewidth]{figure_06_qr_core_equivalence.png}
\caption{QR-core spectral equivalence under the stated orthonormality hypotheses. The diagram pairs the exact factor identity with the archived maximum absolute and relative \(\ell_2\) singular-spectrum discrepancies for a 128-row factor with 180 old and 37 appended columns. Recorded audit metrics replace the overlaid singular-value curves, whose raw vectors were not archived.}
\label{fig:qr-core-equivalence}
\end{figure}

\section{Range-core spectral compression}\label{sec:final-21}

Suppose an orthonormal basis \(Q_r\) captures the factor range:
\[
\operatorname{range}(A)\subseteq\operatorname{range}(Q_r).
\]
Then
\[
A=Q_rB,
\qquad
B=Q_r^\top A.
\]

\begin{theorem}[Exact range core]\label{claim:final-21-1}
The nonzero eigenvalues of
\[
K=AA^\top
\]
are exactly the eigenvalues of
\[
BB^\top.
\]
\end{theorem}

Thus the computational pipeline can terminate its data-streaming stage after
the range is captured:
\[
\boxed{
\text{streamed data}
\to
\text{small spectral core}
\to
\text{tail evaluation}.
}
\]

\section{From phase pruning to certified quadratic-form tails}\label{sec:final-22}

The previous sections certify the effect of phase pruning on Krylov
propagation.  A SKAT-like test ultimately depends on the spectrum itself, so
one more bridge is needed:
\[
\text{phase pruning}
\longrightarrow
\text{operator error}
\longrightarrow
\text{spectral error}
\longrightarrow
\text{tail-probability error}.
\]

\subsection{Block-norm operator domination}

Let
\[
E=H-H_0
\]
be the self-adjoint error introduced by phase pruning.  Using the same
orthogonal phase projectors \(P_h\), define
\[
\mathcal C_{hg}
=
\|P_hEP_g\|_2.
\]

\begin{theorem}[Block-norm operator certificate]\label{claim:final-22-1}
\[
\boxed{
\|E\|_2
\le
\|\mathcal C\|_2.
}
\]
\end{theorem}

\begin{proof}
For any vector \(x\), let
\[
r_g=\|P_gx\|_2.
\]
Then
\[
\|P_hEx\|_2
\le
\sum_g
\|P_hEP_g\|_2\,\|P_gx\|_2
=
(\mathcal C r)_h.
\]
Orthogonality of the phase decomposition gives
\[
\|Ex\|_2^2
=
\sum_h\|P_hEx\|_2^2
\le
\|\mathcal C r\|_2^2
\le
\|\mathcal C\|_2^2\|r\|_2^2
=
\|\mathcal C\|_2^2\|x\|_2^2.
\]
Taking the supremum over unit vectors proves the claim.
\end{proof}

The important computational point is that \(\mathcal C\) is only a
phase-by-phase matrix.  The same block norms already used to construct the
polyphase scheduler therefore provide a global operator certificate.

\subsection{Weyl enclosure of the quadratic-form coefficients}

Assume
\[
H\succeq0
\]
and let the eigenvalues of \(H\) and \(H_0\), in decreasing order, be
\[
\lambda_i
\qquad\text{and}\qquad
\mu_i.
\]
If
\[
\|H-H_0\|_2\le\varepsilon,
\]
Weyl's inequality gives
\[
|\lambda_i-\mu_i|\le\varepsilon.
\]

\begin{corollary}[Positive spectral interval]\label{claim:final-22-2}
For every \(i\),
\[
\boxed{
\ell_i
=
\max\{0,\mu_i-\varepsilon\}
\le
\lambda_i
\le
u_i
=
\max\{0,\mu_i+\varepsilon\}.
}
\]
\end{corollary}

This remains valid even if symmetric phase pruning makes \(H_0\) slightly
indefinite.

\subsection{Certified generalized-\texorpdfstring{\(\chi^2\)}{chi-square} tail bracket}

For independent standard normal variables \(Z_i\), define
\[
Q_-=\sum_i\ell_iZ_i^2,
\qquad
Q=\sum_i\lambda_iZ_i^2,
\qquad
Q_+=\sum_i u_iZ_i^2.
\]

\begin{theorem}[Stochastic tail enclosure]\label{claim:final-22-3}
For every threshold \(q\),
\[
\boxed{
\Pr(Q_-\ge q)
\le
\Pr(Q\ge q)
\le
\Pr(Q_+\ge q).
}
\]
\end{theorem}

\begin{proof}
Couple all three quadratic forms using the same normal variables \(Z_i\).
Since
\[
0\le\ell_i\le\lambda_i\le u_i
\]
for every \(i\),
\[
Q_-\le Q\le Q_+
\]
almost surely.  Taking upper-tail events yields the result.
\end{proof}

The two endpoint probabilities can be evaluated by the same K-COS machinery
used for the unpruned spectrum.  Its numerical error certificate is separate
from the spectral enclosure above.

\subsection{A dual-budget pruning rule}

The local Krylov certificate from the previous section and the global spectral
certificate answer different questions.  The former asks whether omitted
couplings matter to the current Krylov query; the latter asks whether they can
move the spectrum enough to matter to the final test.

A production schedule should therefore enforce
\[
\boxed{
\sum_{e\in\mathcal O}\beta_e^{(k)}
\le
\tau_K,
\qquad
\|\mathcal C_{\mathcal O}\|_2
\le
\tau_\lambda.
}
\]

The implementation uses a conservative greedy rule: phase pairs are considered
in increasing query-risk order, and a pair is omitted only when adding it
preserves both inequalities.  The rule is certified but is not claimed to be
globally optimal under the two simultaneous constraints.

\begin{figure}[H]
\centering
\includegraphics[width=.84\linewidth]{figure_14_dual_budget_tail_bridge.png}
\caption{Dual-budget phase pruning.  The query-local budget controls Krylov
propagation while the operator budget controls final spectral and tail
uncertainty.  A loose operator budget can admit a strong coupling and widen
the tail interval by several decades even when the Krylov-local budget still
passes.}
\label{fig:dual-budget-tail}
\end{figure}

In the deterministic ten-phase experiment, fixing the query budget at
\(0.05\) and the operator budget at \(0.005\) removed \(23\) of \(45\)
undirected phase pairs.  At a true probability of \(10^{-8}\), the certified
tail interval was
\[
10^{-8.11143}
\le
p
\le
10^{-7.89036},
\]
a width of about \(0.221\) decades.  With an operator budget of \(0.0075\),
\(34\) pairs were removed and the interval widened to about \(0.320\)
decades.  By contrast, allowing a strong coupling to be removed under a very
loose operator budget produced a multi-decade interval.

\subsection{Certify after prune}

There is a second route that does not use the pruned operator as the final
spectral approximation.  Let \(U\) be any orthonormal sample-space basis,
possibly generated cheaply from a pruned or approximate Krylov process, and
define
\[
B=U^\top A,
\qquad
F=(I-UU^\top)A,
\qquad
\delta=\|F\|_2.
\]

\begin{theorem}[Range-core singular-value certificate]\label{claim:final-22-4}
Let the singular values of \(A\) be \(\sigma_i(A)\), and let
\(\widetilde\sigma_i(B)\) denote the singular values of \(B\) padded by zeros
to the ambient spectral dimension.  Then
\[
\boxed{
\left|
\sigma_i(A)
-
\widetilde\sigma_i(B)
\right|
\le
\delta.
}
\]
Consequently,
\[
\max\{0,\widetilde\sigma_i(B)-\delta\}^2
\le
\lambda_i(AA^\top)
\le
\left(
\widetilde\sigma_i(B)+\delta
\right)^2.
\]
\end{theorem}

\begin{proof}
Because \(U\) has orthonormal columns,
\[
A=UB+F
\]
and
\[
\|A-UB\|_2=\|F\|_2=\delta.
\]
The nonzero singular values of \(UB\) are exactly those of \(B\).
The standard singular-value perturbation inequality gives the stated bound.
Squaring the nonnegative endpoints gives the eigenvalue interval.
\end{proof}

This provides a useful architecture:
\[
\boxed{
\text{pruned operator for basis discovery}
\to
\text{one exact streamed factor check}
\to
\text{certified spectral core}.
}
\]

\subsection{A retained negative result}

The preceding theorem is exact, but a correct interval need not be useful.

In a full-rank \(50\times50\) stress test, a pruned Krylov basis of dimensions
\(12\) through \(48\) left range residuals between roughly \(0.87\) and
\(0.99\) times the leading singular value.  The corresponding \(10^{-8}\)
tail intervals were therefore vacuous or thousands of decades wide.  At
dimension \(50\), the range residual fell to approximately
\(1.2\times10^{-15}\), and the tail interval collapsed to the exact value.

\begin{figure}[H]
\centering
\includegraphics[width=.84\linewidth]{figure_15_range_core_certification.png}
\caption{Certify-after-prune range residual.  The range-core theorem is valid
at every depth, but a deep-tail full-spectrum guarantee can remain useless
until the range residual is genuinely small.  This negative result is retained
as a stopping-rule lesson rather than hidden as an implementation failure.}
\label{fig:range-core-cert}
\end{figure}

The practical lesson is that range-core certification is a stopping test, not
a promise that a severely incomplete basis automatically yields a useful
full-spectrum tail bound.


\section{Trace-aware residual certification and exact completion}\label{sec:final-23}

The spectral bridge in the preceding section used an operator-norm residual as
if every unresolved coefficient could move independently by the same amount.
That is rigorous but unnecessarily pessimistic once a nontrivial range basis
has already been captured.  The residual has considerably more structure.

Let \(U\in\mathbb R^{n\times r}\) have orthonormal columns and define
\[
B=U^\top A,
\qquad
F=(I-UU^\top)A.
\]
Then
\[
A=UB+F,
\qquad
U^\top F=0.
\]

\begin{theorem}[Exact core--residual decomposition]\label{claim:final-23-1}
\[
\boxed{
A^\top A
=
B^\top B+F^\top F.
}
\]
In particular the unresolved spectral term
\[
R=F^\top F
\]
is positive semidefinite.
\end{theorem}

\begin{proof}
Expand \((UB+F)^\top(UB+F)\).  The two cross terms vanish because
\(U^\top F=0\), while \(U^\top U=I\).
\end{proof}

This changes the certification problem.  Instead of an arbitrary Hermitian
error, the missing term is a positive residual with measurable spectral
moments.  Define
\[
t=\operatorname{tr}R,
\qquad
v=\operatorname{tr}(R^2),
\qquad
\rho=\|R\|_2.
\]
Then
\[
\boxed{
 t=\|F\|_F^2,
 \qquad
 v=\|FF^\top\|_F^2,
 \qquad
 \rho=\|F\|_2^2.
}
\]
The first quantity is the exact missing mean of the Gaussian quadratic form;
the second records how concentrated the residual spectrum is; the third is its
largest possible coefficient.

\subsection{A trace--Frobenius--cap MGF envelope}

Set
\[
C=B^\top B,
\qquad
H=C+R,
\qquad
c=\lambda_{\max}(C).
\]
For \(0<2\theta(c+\rho)<1\), put
\[
\alpha=(1-2\theta c)^{-1}.
\]

\begin{theorem}[Trace--Frobenius--cap log-MGF certificate]\label{claim:final-23-2}
For the Gaussian quadratic form with coefficient matrix \(H\),
\[
K_H(\theta)
=-\frac12\log\det(I-2\theta H)
\]
satisfies
\[
\boxed{
K_H(\theta)
\le
K_C(\theta)
+
\theta\alpha t
+
\frac{\theta^2\alpha^2v}
     {1-2\theta\alpha\rho}.
}
\]
\end{theorem}

\begin{proof}
Factor the determinant through the known core:
\[
I-2\theta(C+R)
=
(I-2\theta C)^{1/2}
(I-2\theta T)
(I-2\theta C)^{1/2},
\]
where
\[
T=(I-2\theta C)^{-1/2}
R
(I-2\theta C)^{-1/2}.
\]
Since
\[
\|(I-2\theta C)^{-1/2}\|_2^2=\alpha,
\]
we have
\[
\operatorname{tr}T\le\alpha t,
\qquad
\operatorname{tr}(T^2)\le\alpha^2v,
\qquad
\|T\|_2\le\alpha\rho.
\]
For every eigenvalue \(x\) of \(T\),
\[
-\frac12\log(1-2\theta x)
\le
\theta x
+
\frac{\theta^2x^2}
     {1-2\theta\alpha\rho}.
\]
The scalar inequality follows from the power series of \(-\log(1-u)\) and a
geometric bound on all terms of degree at least two.  Summing over the
eigenvalues of \(T\) proves the result.
\end{proof}

\begin{corollary}[Trace-aware Chernoff upper tail]\label{claim:final-23-3}
For \(Q=Z^\top HZ\) with \(Z\sim N(0,I)\),
\[
\boxed{
\Pr(Q\ge q)
\le
\inf_{0<2\theta(c+\rho)<1}
\exp\{\widehat K_H(\theta)-\theta q\},
}
\]
where \(\widehat K_H\) is the preceding MGF upper bound.
\end{corollary}

A complementary lower endpoint is immediate from
\[
H=C+R\succeq C.
\]
Using the same normal vector,
\[
Z^\top HZ\ge Z^\top CZ
\]
almost surely, and therefore
\[
\boxed{
\Pr(Z^\top CZ\ge q)
\le
\Pr(Z^\top HZ\ge q).
}
\]
The lower endpoint can be evaluated by K-COS on the captured core spectrum,
while the upper endpoint is obtained from the optimized trace-aware Chernoff
bound.

\subsection{Finite stress test}

The previous coefficient-wise Weyl interval was intentionally conservative.
On the same full-rank, 50-dimensional stress test used in the preceding
section, its \(10^{-8}\) tail interval remained thousands of decades wide
until essentially the complete range had been captured.  Preserving the exact
residual trace and higher residual moments changes that behavior substantially.

\begin{table}[H]
\centering
\caption{Tail-bracket width at a true probability of \(10^{-8}\).  The
trace/moment certificate is still conservative, but it is orders of magnitude
narrower than coefficient-wise Weyl inflation before full range capture.}
\begin{tabular}{rrrr}
\toprule
Krylov depth & basis dimension & v0.7 Weyl width & v0.8 moment width\\
\midrule
4 & 24 & 90,097 & 11.86\\
5 & 30 & 16,511 & 9.76\\
6 & 36 & 3,747 & 7.67\\
7 & 42 & 2,275 & 5.70\\
8 & 48 & 1,165 & 3.39\\
\bottomrule
\end{tabular}
\end{table}

\begin{figure}[H]
\centering
\includegraphics[width=.84\linewidth]{figure_16_trace_moment_vs_weyl.png}
\caption{Incomplete-range tail certification before exact completion.  The
trace/moment bound preserves the amount and concentration of unresolved
spectral mass instead of shifting every coefficient by one worst-case radius.}
\label{fig:trace-moment-vs-weyl}
\end{figure}

A separate random implementation audit used 120 random factor/range pairs and
four admissible positive tilts per pair, for 480 MGF comparisons.  Both the
trace/cap and trace--Frobenius--cap envelopes dominated the exact log MGF in
every tested case.  The trace--Frobenius--cap form was at least as tight as the
trace/cap form in approximately 83.5\% of those comparisons.  This finite
sweep audits the implementation; the theorem above supplies the universal
argument.

\subsection{Residual rank from moments}

The same invariants determine when it is better to stop approximating the
missing spectrum and complete it exactly.  Let
\[
s=\operatorname{rank}F=\operatorname{rank}R.
\]
If the nonzero eigenvalues of \(R\) are \(r_1,\ldots,r_s\), then
\[
t=\sum_i r_i,
\qquad
v=\sum_i r_i^2,
\qquad
r_i\le\rho.
\]

\begin{theorem}[Residual-rank bracket]\label{claim:final-23-4}
\[
\boxed{
\left\lceil
\max\left\{
\frac{t}{\rho},
\frac{t^2}{v}
\right\}
\right\rceil
\le
\operatorname{rank}F
\le
n-r.
}
\]
The zero-residual case is interpreted separately as rank zero.
\end{theorem}

\begin{proof}
The inequality \(t\le s\rho\) gives \(s\ge t/\rho\).  Cauchy--Schwarz gives
\[
t^2
=
\left(\sum_{i=1}^s r_i\right)^2
\le
s\sum_{i=1}^s r_i^2
=sv,
\]
so \(s\ge t^2/v\).  Finally, every column of \(F\) lies in
\(\operatorname{range}(U)^\perp\), whose dimension is \(n-r\).
\end{proof}

Whenever the integer lower bound equals \(n-r\), the missing rank is certified
exactly without computing a residual SVD.

\begin{figure}[H]
\centering
\includegraphics[width=.84\linewidth]{figure_17_residual_rank_completion.png}
\caption{Residual-rank certification in the full-rank stress test.  From depth
3 onward the moment lower bound meets the orthogonal-complement upper bound.
At depth 8 exactly two residual directions remain.}
\label{fig:residual-rank-completion}
\end{figure}

\subsection{Exact residual completion}

Suppose a compact factorization of the residual is
\[
F=WD,
\]
where \(W\) has orthonormal columns.  Since \(F\) lies in the orthogonal
complement of \(U\),
\[
U^\top W=0.
\]

\begin{theorem}[Exact residual completion]\label{claim:final-23-5}
\[
\boxed{
A
=
[U\;W]
\begin{bmatrix}
B\\D
\end{bmatrix}.
}
\]
Consequently, the nonzero singular values of \(A\) are exactly the singular
values of the stacked completion core
\[
\mathcal C_{\rm complete}
=
\begin{bmatrix}
B\\D
\end{bmatrix}.
\]
\end{theorem}

This identity admits a column-streamed implementation.  For each residual
column block, project against the residual basis accumulated so far, append
only genuinely new directions through rank-revealing QR or SVD, and extend the
coordinate matrices of earlier blocks with zero rows.  New residual directions
are orthogonal to every earlier residual block by construction, so no earlier
data block must be revisited.

In the deterministic stress test, exact residual completion reproduced the
full singular spectrum at relative errors between approximately
\(2\times10^{-16}\) and \(6\times10^{-16}\) for all tested incomplete bases.
The resulting \(10^{-8}\) tail differed from the full-spectrum reference by
less than \(10^{-14}\) decades.  At depth 8 the candidate range had dimension
48, the missing rank was certified as exactly two, and appending those two
residual directions closed the spectrum without another approximate Krylov
stage.

\subsection{Hybrid stopping rule}

The resulting production logic is not ``approximate forever.''  It is a
certified handoff between two regimes:
\[
\boxed{
\begin{array}{c}
\text{streamed / pruned Krylov}\\[2mm]
\downarrow\\[-1mm]
\text{residual moments }(t,v,\rho)\\[2mm]
\downarrow\\[-1mm]
\begin{cases}
\text{continue Krylov}, & \text{residual range still large},\\
\text{exact streamed residual completion}, & \text{completion rank affordable}.
\end{cases}
\end{array}
}
\]

The trace-aware tail certificate remains useful as a rigorous intermediate
bound and as a stopping diagnostic.  Once only a small number of residual
directions remain, exact completion is mathematically preferable to forcing a
worst-case tail envelope to become artificially sharp.

\section{Rare-variant kernel case study}\label{sec:final-24}

The preceding sections establish operator identities and certified numerical
bounds.  This section tests the same machinery on the project's existing
rare-variant SKAT simulation layer, including explicit LD structure,
covariates, per-variant standardization, rare-variant weighting, and the
phenotype-space score statistic.

\subsection{Simulation provenance and implementation correction}

Two established genotype generators were used.

The first is a latent-Gaussian block-LD generator.  Within a genomic block,
alleles are obtained by thresholding correlated Gaussian latent variables at
variant-specific minor-allele frequencies.

The second is a copy-LD generator in which haplotype states persist along a
block with high probability and are occasionally redrawn.  During the v0.9
audit, a reproducibility defect was found in an earlier simulator implementation: the genotype
matrix was allocated with \verb|np.empty| and then updated using \verb|+=|.
Because the initial contents of \verb|np.empty| are unspecified, those old
copy-LD benchmark values are not reproducible evidence.  The v0.9 generator
uses an explicitly zero-initialized matrix before haplotype accumulation.
All copy-LD results reported below use the corrected implementation.

\subsection{Tested LD regimes}

Four deterministic scenarios were evaluated.

\begin{table}[H]
\centering
\caption{Realistic rare-variant simulation scenarios.  Adjacent $r^2$ is
computed within the simulated LD blocks after variant standardization.}
\begin{tabular}{lrrrrr}
\toprule
scenario & $n$ & $m$ & rank & median adjacent $r^2$ & effective rank\\
\midrule
latent LD, moderate & 600 & 3000 & 596 & 0.0069 & 441.9\\
latent LD, strong & 600 & 3000 & 596 & 0.1140 & 240.3\\
copy LD, strong & 600 & 3000 & 596 & 0.8298 & 200.1\\
copy LD, large & 1000 & 5000 & 996 & 0.8179 & 353.0\\
\bottomrule
\end{tabular}
\end{table}

The corresponding spectra therefore span a substantial range of concentration:
stronger LD reduces effective rank even when the nominal quadratic-form rank
remains close to the sample residual dimension.

\begin{figure}[H]
\centering
\includegraphics[width=.84\linewidth]{figure_18_realistic_ld_effective_rank.png}
\caption{LD structure changes the effective spectral rank.  The corrected
copy-LD models create much stronger local dependence than the latent-Gaussian
moderate-LD case, while all remain within the same streamed SKAT pipeline.}
\label{fig:realistic-ld-rank}
\end{figure}

\subsection{Streaming directly from raw genotype blocks}

The raw genotype matrix was streamed in resident variant blocks of widths
31, 127, and 509.  Within each block the implementation performs

\[
G_b
\to
G_b-\bar G_b
\to
\text{variance scaling}
\to
M_X G_b
\to
W_b^{1/2}
\to
A_bA_b^\top,
\]

without materializing the full weighted factor.

Across all four LD scenarios and all three resident block sizes, the largest
relative discrepancy against the fully materialized kernel was

\[
\boxed{6.61\times10^{-16}}.
\]

\begin{figure}[H]
\centering
\includegraphics[width=.84\linewidth]{figure_19_realistic_streaming_invariance.png}
\caption{Raw-genotype streaming reproduces the materialized SKAT kernel across
multiple LD regimes and block sizes.  The observed discrepancies are at
floating-point summation scale.}
\label{fig:realistic-streaming}
\end{figure}

This test is stronger than applying the streamed identity to a precomputed
factor because the centering, variance scaling, covariate projection, and
rare-variant weights are all performed inside the block loop.

\subsection{Phenotype-space equivalence}

For each scenario, 100 Gaussian null residual draws were evaluated in three
ways:

\begin{enumerate}
\item directly through the factor,
\[
Q=\frac12\|A^\top e\|_2^2;
\]
\item through the sample kernel,
\[
Q=\frac12 e^\top Ke;
\]
\item through the eigencoordinates of $K$.
\end{enumerate}

The worst relative discrepancy across all draws and all scenarios was

\[
\boxed{8.32\times10^{-16}}.
\]

Thus the simulated phenotype path, the streamed kernel, and the generalized
$\chi^2$ spectral representation agree to floating-point precision.

\subsection{Deep tails on realistic LD spectra}

For every scenario, tail queries were generated at tilt fractions
0.90, 0.97, and 0.99 of the positive singularity.  K-COS was compared with a
higher-resolution full-spectrum tilted reference.

The large corrected copy-LD case was then pushed further, to fractions
0.995, 0.997, 0.999, and 0.9995.  The deepest tested reference value was

\[
\boxed{\log_{10}p=-457.973916}.
\]

At that point the displayed K-COS and high-resolution full-spectrum values
were identical.  Across the complete realistic-LD battery, the largest
absolute discrepancy was

\[
\boxed{3.21\times10^{-12}\ \text{decades}}.
\]

\begin{figure}[H]
\centering
\includegraphics[width=.84\linewidth]{figure_20_realistic_deep_tail_accuracy.png}
\caption{K-COS versus the high-resolution full-spectrum tilted reference on
realistic LD spectra.  Accuracy remains at roughly machine-scale in log
probability through the deepest corrected copy-LD query, near $10^{-458}$.}
\label{fig:realistic-deep-tail}
\end{figure}

For the $n=1000$, $m=5000$ corrected copy-LD spectrum, selected results are

\begin{center}
\begin{tabular}{rrrr}
\toprule
tilt fraction & $\log_{10}p$ & K-COS $K$ & absolute error (decades)\\
\midrule
0.9950 & -65.950669 & 2048 & $4.69\times10^{-13}$\\
0.9970 & -95.305034 & 2048 & $3.70\times10^{-14}$\\
0.9990 & -240.602047 & 8192 & $4.94\times10^{-14}$\\
0.9995 & -457.973916 & 16384 & below stored reference resolution\\
\bottomrule
\end{tabular}
\end{center}

\subsection{Interpretation and boundary}

This simulation block does not add a new algebraic theorem.  It supplies a
more demanding validation layer for the theorem chain already proved.
Specifically, it shows that the exact streaming identities and the K-COS
deep-tail machinery remain numerically stable under substantially different
rare-variant LD structures and through the complete raw-genotype preprocessing
path.

The evidence remains synthetic.  These generators do not reproduce every
population-genetic feature of a sequencing cohort, and this block is not an
external-cohort replication.  It also addresses the Gaussian linear-model null;
finite-sample Bernoulli or logistic calibration remains a separate statistical
problem.


\section{General production API and deep-tail arithmetic hardening}\label{sec:final-25}

The preceding mathematics is exposed through a domain-neutral engine rather
than through a SKAT-only entry point.  The central positive engine accepts:

\begin{itemize}
\item an explicit nonnegative spectrum;
\item a positive semidefinite matrix;
\item a dense factor \(A\);
\item a replayable streamed column-factor source;
\item or a covariance/quadratic-matrix pair
\((\Sigma,M)\).
\end{itemize}

A signed central engine accepts a signed spectrum, a symmetric operator, or
\[
(\Sigma,M),\qquad M=M^\top,
\]
and performs the exact central-indefinite reduction proved in Section~\ref{sec:gaussian-core}.  A
separate noncentral positive engine accepts
\[
(\mu,\Sigma,M),\qquad M\succeq0,
\]
and performs the corresponding noncentral reduction.

\subsection{Log-domain packet error accumulation}

At extreme tail depths the analytic packet certificate contains sums of the
form
\[
\frac{2}{U}
\sum_k
|I_k|
\left(
e^{\varepsilon_k}-1
\right).
\]
Direct evaluation can overflow even when the candidate is destined to be
rejected by the adaptive planner.  The hardened implementation instead
computes
\[
\log\left(
\frac{2}{U}
\sum_k
|I_k|
\left(
e^{\varepsilon_k}-1
\right)
\right)
\]
by log-sum-exp arithmetic, and compares this quantity with the logarithm of
the requested tolerance before converting a final accepted bound back to
float64.

This change does not weaken the mathematical bound; it changes only its
floating-point representation.

\subsection{Reference convergence is part of verification}

An independently evaluated reference is useful only if its own numerical
resolution has converged.  In a rank-350 deep-tail audit, an initial fixed-grid
reference appeared to disagree with the adaptive answer.  Increasing the
reference COS grid gave discrepancies of approximately
\[
1.88\times10^{-2},
\quad
2.6\times10^{-3}\text{--}4.0\times10^{-3},
\quad
1.39\times10^{-4},
\quad
2.60\times10^{-8}
\]
decades as the grid was refined through
\[
K=4096,\ 8192,\ 16384,\ 32768.
\]

\begin{figure}[H]
\centering
\includegraphics[width=.82\linewidth]{figure_21_reference_convergence.png}
\caption{Reference-convergence audit.  The adaptive certified answer remains
stable while an initially under-resolved fixed-grid reference converges toward
it.  A single nominally large reference grid is therefore not treated as
authoritative without its own refinement check.}
\label{fig:reference-convergence}
\end{figure}

The production verification rule is consequently two-sided: the adaptive
solver must satisfy its own certificate, and any external or fixed-grid
reference used for comparison must demonstrate numerical convergence.


\section{Centered spectral packets and K-COS}\label{sec:packets}
Matrix storage packets operate before eigendecomposition.  The spectral
packets here operate after the eigenvalues are available and compress the
\emph{transform}, not the probability law into a moment-matched surrogate.
The same operator can therefore use storage blocks, structural phases, range
compression, and spectral packets simultaneously.

The packet results in this section are mathematical approximation theorems.
Their binary64 implementation reports a separate analytic ledger.  The
directed-rounding positive backend developed earlier in this manuscript currently evaluates the
supplied positive spectrum directly; it does not automatically validate a
binary64 packet approximation.  A combined directed packet implementation
must enclose packet centers, moments, remainders, and transform arithmetic.

\begin{figure}[H]
\centering
\includegraphics[width=.86\linewidth]{figure_07_spectral_packets.png}
\caption{A second decomposition layer: centered spectral moments compress the
complex transform after the operator problem has produced a spectrum.  This
schematic uses the definitions for nonempty packets and preserves the rounded
packet-center and second-moment labels from the archived illustration,
independently of physical storage blocks. It does not reconstruct the
unavailable original eigenvalue vector.}
\label{fig:spectral-packets}
\end{figure}

\subsection{Quadratic-form transform and saddlepoint tilt}

Let
\begin{equation}
Q=\sum_{j=1}^{r}\lambda_j Z_j^2,
\qquad \lambda_j>0.
\label{eq:Q}
\end{equation}
The moment generating function exists for
\[
\Re s < \frac{1}{2\lambda_{\max}},
\qquad
\lambda_{\max}=\max_j\lambda_j,
\]
and its cumulant generating function is
\begin{equation}
\kappa(s)=\log \E e^{sQ}
=-\frac12\sum_{j=1}^{r}\log(1-2s\lambda_j).
\label{eq:cgf}
\end{equation}
For real $\theta<1/(2\lambda_{\max})$,
\begin{equation}
\kappa^{(m)}(\theta)
=2^{m-1}(m-1)!\sum_{j=1}^{r}
\frac{\lambda_j^m}{(1-2\theta\lambda_j)^m},
\qquad m\ge1.
\label{eq:derivatives}
\end{equation}
At $\theta=0$, this recovers the standard cumulant identity
\[
\operatorname{cum}_m(Q)=2^{m-1}(m-1)!\sum_j\lambda_j^m.
\]

\begin{theorem}[Unique upper-tail saddlepoint]
\label{thm:saddle}
If $q>\E Q=\sum_j\lambda_j$, there exists a unique
\[
\theta_*=\theta(q)\in\left(0,\frac{1}{2\lambda_{\max}}\right)
\]
satisfying
\begin{equation}
\kappa'(\theta_*)=q.
\label{eq:saddle}
\end{equation}
\end{theorem}

\begin{proof}
By \eqref{eq:derivatives},
\[
\kappa'(\theta)=\sum_j\frac{\lambda_j}{1-2\theta\lambda_j}
\]
is continuous and strictly increasing on the stated interval. It equals $\E Q$ at $\theta=0$ and diverges as $\theta\uparrow1/(2\lambda_{\max})$. The intermediate value theorem and strict monotonicity give existence and uniqueness.
\end{proof}

Define the exponentially tilted law by
\begin{equation}
f_{\theta}(x)=e^{\theta x-\kappa(\theta)}f(x).
\label{eq:tilt-density}
\end{equation}
Its characteristic function is
\begin{equation}
\phi_{\theta}(t)
=\exp\{\kappa(\theta+\ii t)-\kappa(\theta)\},
\label{eq:tilt-cf}
\end{equation}
and $\E_{\theta}Q=\kappa'(\theta)$. Therefore the choice \eqref{eq:saddle} moves the query $q$ to the mean of the tilted law.

For any $\theta>0$ in the transform domain,
\begin{equation}
\Pp(Q\ge q)
=e^{\kappa(\theta)-\theta q}
\E_{\theta}\!\left[e^{-\theta(Q-q)}\mathbf 1_{\{Q\ge q\}}\right].
\label{eq:tilt-tail}
\end{equation}
This identity is classical exponential tilting; K-COS uses it to avoid subtracting a CDF value numerically close to one.

\subsection{Centered spectral packets}

Let a nonempty packet $P\subset\{1,\ldots,r\}$ contain $n_P\ge1$ eigenvalues. Define its center
\[
\mu_P=\frac{1}{n_P}\sum_{j\in P}\lambda_j,
\qquad
\delta_j=\lambda_j-\mu_P,
\]
and centered moments
\begin{equation}
M_{P,m}=\sum_{j\in P}\delta_j^m,
\qquad m\ge0.
\label{eq:centeredmoments}
\end{equation}
Thus $M_{P,0}=n_P$ and $M_{P,1}=0$.

\begin{theorem}[Centered packet expansion]
\label{thm:packet}
Fix $s\in\C$ with $\Re s<1/(2\lambda_{\max})$, use the principal logarithm, and define
\begin{equation}
\alpha_P(s)=\frac{2s}{1-2s\mu_P},
\qquad
\rho_P(s)=|\alpha_P(s)|\max_{j\in P}|\delta_j|.
\label{eq:alpha-rho}
\end{equation}
For a nonnegative integer expansion order $p$, if $\rho_P(s)<1$, then the packet contribution to \eqref{eq:cgf} is
\begin{align}
\kappa_P(s)
&=-\frac{n_P}{2}\log(1-2s\mu_P)
+\frac12\sum_{m=1}^{p}
\frac{\alpha_P(s)^m}{m}M_{P,m}
+R_{P,p}(s),
\label{eq:packet-expansion}
\end{align}
where
\begin{equation}
|R_{P,p}(s)|
\le
\frac{n_P}{2}
\frac{\rho_P(s)^{p+1}}{(p+1)(1-\rho_P(s))}.
\label{eq:packet-remainder}
\end{equation}
\end{theorem}

\begin{proof}
For $j\in P$,
\[
1-2s\lambda_j
=(1-2s\mu_P)(1-\alpha_P(s)\delta_j).
\]
The transform-domain condition gives positive real part for each
\(1-2s\lambda_j\) and for \(1-2s\mu_P\).  The disk condition gives
positive real part for \(1-\alpha_P\delta_j\).  The two factor arguments
therefore sum inside \((-\pi,\pi)\), so their principal logarithms add
without a branch ambiguity.  Hence
\[
-\frac12\log(1-2s\lambda_j)
=-\frac12\log(1-2s\mu_P)
-\frac12\log(1-\alpha_P\delta_j).
\]
When $|\alpha_P\delta_j|<1$,
\[
-\log(1-z)=\sum_{m\ge1}\frac{z^m}{m}.
\]
Summing over $j\in P$ gives \eqref{eq:packet-expansion}. The remainder satisfies
\[
\left|\sum_{m\ge p+1}\frac{z^m}{m}\right|
\le
\sum_{m\ge p+1}\frac{|z|^m}{p+1}
=\frac{|z|^{p+1}}{(p+1)(1-|z|)}.
\]
Using $|z|\le\rho_P(s)$ and summing $n_P$ terms yields \eqref{eq:packet-remainder}.
\end{proof}

\begin{remark}[What is compressed]\label{claim:final-26-3}
A packet of $n_P$ eigenvalues is replaced by its center and $p$ centered moments. Large spectral spikes may be retained exactly. If the representation contains $s$ exact eigenvalues and $R$ packets of order $p$, evaluation of the CGF on $N$ frequencies costs $O(N(s+Rp))$ arithmetic operations after compilation, instead of $O(Nr)$ direct eigenvalue-transform evaluations.
\end{remark}

\begin{remark}[Query dependence]\label{claim:final-26-4}
The relevant packet condition is not merely geometric closeness of eigenvalues. It is the complex-grid quantity $\rho_P(\theta_*+\ii u_k)$. Thus a partition suitable for one tail query or one COS grid need not be suitable for another.
\end{remark}

\subsection{Tilted COS tail functional}

Because $Q\ge0$, K-COS uses the interval $[0,b]$.  Throughout the
following functional and packet-error results, assume $b>q\ge0$, $b>0$,
$0<\theta<1/(2\lambda_{\max})$, and an integer mode count $N\ge1$.  Let
\[
u_k=\frac{k\pi}{b},
\qquad k=0,1,\ldots,
\]
and define the exponentially weighted cosine integrals
\begin{equation}
J_k(q,b,\theta)
=\int_q^b e^{-\theta(x-q)}\cos(u_kx)\,dx.
\label{eq:J-def}
\end{equation}
They admit the closed form
\begin{equation}
J_k
=\Re\left[
 e^{\ii u_k q}
 \frac{1-e^{(-\theta+\ii u_k)(b-q)}}{\theta-\ii u_k}
\right].
\label{eq:J-formula}
\end{equation}
The prime on a cosine sum denotes half weight for $k=0$.

If the exact tilted characteristic function is used, the $N$-mode finite-window expression is
\begin{equation}
T_{N,b}
=\frac{2}{b}\sum_{k=0}^{N-1}{}'
J_k\,
\Re\left[e^{\kappa(\theta+\ii u_k)-\theta q}\right].
\label{eq:tail-cos}
\end{equation}
The cancellation of $\kappa(\theta)$ in \eqref{eq:tail-cos} is useful numerically: the likelihood-ratio prefactor from \eqref{eq:tilt-tail} and the normalization in \eqref{eq:tilt-cf} do not need to be formed separately inside each mode.

\subsection{Propagation of packet error}

Let $\widehat\kappa(s_k)$ be the packet approximation at
\[
s_k=\theta+\ii u_k,
\]
and suppose
\begin{equation}
|\widehat\kappa(s_k)-\kappa(s_k)|\le B_k.
\label{eq:Bk}
\end{equation}
For disjoint packets, $B_k$ can be chosen as the sum of the packet bounds \eqref{eq:packet-remainder}.

\begin{theorem}[Finite-grid packet-to-tail bound]
\label{thm:packet-tail}
Let $\widehat T_{N,b}$ be obtained from \eqref{eq:tail-cos} by replacing $\kappa(s_k)$ by $\widehat\kappa(s_k)$. Then
\begin{equation}
|\widehat T_{N,b}-T_{N,b}|
\le \eps_{\mathrm{pkt}},
\label{eq:packet-tail-bound}
\end{equation}
where
\begin{equation}
\eps_{\mathrm{pkt}}
=
\frac{2}{b}
 e^{\kappa(\theta)-\theta q}
\sum_{k=0}^{N-1}{}'
|J_k|\bigl(e^{B_k}-1\bigr).
\label{eq:eps-pkt}
\end{equation}
\end{theorem}

\begin{proof}
Write $\Delta_k=\widehat\kappa(s_k)-\kappa(s_k)$. Then
\[
|e^{\widehat\kappa(s_k)-\theta q}-e^{\kappa(s_k)-\theta q}|
\le
|e^{\kappa(s_k)-\theta q}|\,(e^{|\Delta_k|}-1).
\]
Moreover
\[
\Re\kappa(\theta+\ii u_k)-\kappa(\theta)
=\log|\phi_{\theta}(u_k)|\le0,
\]
so
\[
|e^{\kappa(s_k)-\theta q}|
\le e^{\kappa(\theta)-\theta q}.
\]
Substitute this into the finite cosine sum and apply the triangle inequality.
\end{proof}

\subsection{Positive support, truncation, and cosine periodization}

Using full characteristic-function coefficients on $[0,b]$ reconstructs the even $2b$-periodization of the tilted density rather than the strictly truncated density. For a nonnegative density supported on $[0,\infty)$, the total mass folded back into $[0,b]$ equals the tilted mass above $b$. This gives a clean error control for the positive quadratic-form setting.


\begin{lemma}[Integrated cosine identity by folding]
\label{lem:folding}
Let \(\nu_\theta\) be the tilted probability measure and let
\(F_b:[0,\infty)\to[0,b]\) fold each interval of length \(2b\) by
reflection.  Then
\(\cos(k\pi F_b(x)/b)=\cos(k\pi x/b)\).
For \(h(x)=e^{-\theta(x-q)}\mathbf1_{[q,b]}(x)\),
\[
\int h\,d(F_{b\#}\nu_\theta)
=\frac2b\sum_{k=0}^{\infty}{}'
J_k\,\Re\phi_\theta(k\pi/b).
\]
The series is absolutely convergent for every positive rank \(r\ge1\).
\end{lemma}
\begin{proof}
A positive weighted chi-square sum has an absolutely continuous law; its
folded measure therefore has no atom at the jump of \(h\).  The even periodic
extension of \(h\) has bounded Fej\'er means, which converge at every continuity
point.  Dominated convergence identifies the integral of those means with
\(\int h\,d(F_{b\#}\nu_\theta)\).  Their cosine coefficients are the Ces\`aro
weighted coefficients \(2J_k/b\), and folding preserves the cosine moments
\(\Re\phi_\theta(k\pi/b)\).  Finally
\(J_k=O(k^{-1})\) and
\(|\phi_\theta(k\pi/b)|=O(k^{-r/2})\), as follows directly from the product
formula below.  Thus the coefficient products are
\(O(k^{-1-r/2})\); their absolute sum converges.  The Ces\`aro sums and ordinary
sum have the same limit, proving the displayed identity.  This argument
justifies the weighted functional without requiring pointwise convergence of
the reconstructed density at every boundary.
\end{proof}

\begin{proposition}[Support and periodization bound]
\label{prop:support}
Let the finite positive spectrum have rank \(r\ge1\), and let
\(b>q\ge0\), \(b>0\), and \(0<\theta<1/(2\lambda_{\max})\).  Let
\[
\tau_{\theta}(b)=\Pp_{\theta}(Q>b).
\]
Then the combined error from (i) replacing the tilted density on $[0,b]$ by its full-coefficient cosine periodization and (ii) omitting the original integral over $(b,\infty)$ is bounded by
\begin{equation}
\eps_{\mathrm{sup}}
\le
e^{\kappa(\theta)-\theta q}\,
\tau_{\theta}(b)
\left(1+e^{-\theta(b-q)}\right).
\label{eq:eps-support}
\end{equation}
\end{proposition}

\begin{proof}
By Lemma~\ref{lem:folding}, the infinite cosine functional integrates
the folded measure.  The part of that measure originating from \((b,\infty)\)
is nonnegative and has total mass $\tau_{\theta}(b)$. On $[q,b]$, the weight $e^{-\theta(x-q)}$ is at most one, so its contribution to the weighted functional is at most $\tau_{\theta}(b)$. For the omitted original tail, $x\ge b$ implies
\[
e^{-\theta(x-q)}\le e^{-\theta(b-q)},
\]
so that term is at most $e^{-\theta(b-q)}\tau_{\theta}(b)$. Multiplication by the likelihood-ratio prefactor in \eqref{eq:tilt-tail} gives \eqref{eq:eps-support}.
\end{proof}

The tilted tail probability itself can be bounded by Chernoff's inequality. For $\eta>\theta$ inside the transform domain,
\begin{equation}
\tau_{\theta}(b)
\le
\exp\left\{
\kappa(\eta)-\kappa(\theta)-(\eta-\theta)b
\right\}.
\label{eq:tilted-chernoff}
\end{equation}
When $b>\kappa'(\theta)$, the minimizing $\eta$ is the unique solution of
\begin{equation}
\kappa'(\eta)=b.
\label{eq:eta}
\end{equation}
For the saddlepoint tilt, $\kappa'(\theta_*)=q<b$, so such an $\eta$ exists.

\subsection{A specialized omitted-mode bound}

The tilted characteristic function has an especially transparent modulus.

\begin{lemma}[Tilted characteristic-function modulus]
\label{lem:modulus}
Define
\begin{equation}
a_j=\frac{2\lambda_j}{1-2\theta\lambda_j}>0.
\label{eq:aj}
\end{equation}
Then
\begin{equation}
|\phi_{\theta}(t)|
=\prod_{j=1}^{r}\left(1+(a_jt)^2\right)^{-1/4}.
\label{eq:phi-modulus}
\end{equation}
Define the logarithmic decay rate
\begin{equation}
\beta(t)
=-\frac{d\log|\phi_{\theta}(t)|}{d\log t}
=\frac12\sum_{j=1}^{r}
\frac{(a_jt)^2}{1+(a_jt)^2}.
\label{eq:beta}
\end{equation}
Then $\beta(t)$ is nondecreasing for $t\ge0$.
\end{lemma}

\begin{proof}
Taking the modulus of each factor in \eqref{eq:tilt-cf} gives \eqref{eq:phi-modulus}. Differentiating its logarithm gives \eqref{eq:beta}. Each summand $x/(1+x)$ with $x=(a_jt)^2$ is nondecreasing in $t$.
\end{proof}

\begin{corollary}[Power envelope from one frequency]
\label{cor:power-envelope}
For $t\ge t_0>0$,
\begin{equation}
|\phi_{\theta}(t)|
\le
|\phi_{\theta}(t_0)|
\left(\frac{t}{t_0}\right)^{-\beta(t_0)}.
\label{eq:power-envelope}
\end{equation}
\end{corollary}

\begin{proof}
By Lemma \ref{lem:modulus},
\[
\frac{d\log|\phi_{\theta}(t)|}{d\log t}=-\beta(t)\le-\beta(t_0)
\]
for $t\ge t_0$. Integrating from $t_0$ to $t$ proves the claim.
\end{proof}

For $k\ge1$, the exact integral \eqref{eq:J-formula} implies
\begin{equation}
|J_k|
\le \frac{2}{u_k}.
\label{eq:J-bound}
\end{equation}

\begin{theorem}[Omitted COS-mode certificate]
\label{thm:cos-tail}
Let $N\ge1$, $b>q\ge0$, and $u_N=N\pi/b$ and $\beta_N=\beta(u_N)$. If $\beta_N>0$, then the contribution of all exact tilted-COS modes $k\ge N$ to the weighted tail functional is bounded by
\begin{equation}
\eps_{\mathrm{COS}}
\le
 e^{\kappa(\theta)-\theta q}
\frac{4|\phi_{\theta}(u_N)|}{\pi}
\left(\frac1N+\frac1{\beta_N}\right).
\label{eq:eps-cos}
\end{equation}
\end{theorem}

\begin{proof}
Using the cosine coefficient magnitude $2|\phi_{\theta}(u_k)|/b$ and \eqref{eq:J-bound}, the omitted series is at most
\[
e^{\kappa(\theta)-\theta q}
\frac{4}{\pi}
\sum_{k=N}^{\infty}\frac{|\phi_{\theta}(u_k)|}{k}.
\]
By Corollary \ref{cor:power-envelope},
\[
|\phi_{\theta}(u_k)|
\le
|\phi_{\theta}(u_N)|\left(\frac{k}{N}\right)^{-\beta_N}.
\]
Therefore
\[
\sum_{k=N}^{\infty}\frac{|\phi_{\theta}(u_k)|}{k}
\le
|\phi_{\theta}(u_N)|N^{\beta_N}
\sum_{k=N}^{\infty}k^{-\beta_N-1}.
\]
Finally,
\[
\sum_{k=N}^{\infty}k^{-\beta_N-1}
\le N^{-\beta_N-1}+\int_N^{\infty}x^{-\beta_N-1}\,dx
=N^{-\beta_N-1}+\frac{N^{-\beta_N}}{\beta_N},
\]
which yields \eqref{eq:eps-cos}.
\end{proof}

\subsection{Combined analytic certificate}

\begin{theorem}[K-COS analytic error certificate]
\label{thm:combined}
Let $p=\Pp(Q\ge q)$, and let $\widehat p$ be the finite K-COS approximation formed with
\begin{enumerate}[label=(\roman*)]
\item an upper interval endpoint $b>q$,
\item an integer $N\ge1$ of cosine modes,
\item a packet approximation satisfying \eqref{eq:Bk} on the finite grid.
\end{enumerate}
Under the positive-spectrum and parameter hypotheses of Proposition \ref{prop:support},
\begin{equation}
|p-\widehat p|
\le
\eps_{\mathrm{an}}
:=\eps_{\mathrm{pkt}}+\eps_{\mathrm{sup}}+\eps_{\mathrm{COS}}.
\label{eq:combined}
\end{equation}
If $\eps_{\mathrm{an}}<|\widehat p|$, then
\begin{equation}
\frac{|p-\widehat p|}{|p|}
\le
\frac{\eps_{\mathrm{an}}}
{|\widehat p|-\eps_{\mathrm{an}}}.
\label{eq:relative-certificate}
\end{equation}
\end{theorem}

\begin{proof}
Decompose the error into packet error on the retained grid, support and periodization error, and the exact omitted cosine series. The first is bounded by Theorem \ref{thm:packet-tail}, the second by Proposition \ref{prop:support}, and the third by Theorem \ref{thm:cos-tail}. This proves \eqref{eq:combined}. Also
\[
|p|\ge |\widehat p|-|p-\widehat p|
\ge |\widehat p|-\eps_{\mathrm{an}},
\]
which gives \eqref{eq:relative-certificate}.
\end{proof}

\begin{remark}[Certificate boundary]\label{claim:final-26-12}
Theorem \ref{thm:combined} is an \emph{analytic approximation certificate}. It does not enclose IEEE floating-point rounding in complex logarithms, exponentials, root finding, or summation. A machine-level certificate would require directed rounding, interval arithmetic, or a comparable validated-numerics layer.
\end{remark}

\subsection{Finite refinability}

The combined error terms correspond to three distinct obstructions. Packet error is reduced by splitting spectral packets or raising packet order; support error is reduced by increasing $b$; cosine-series error is reduced by increasing $N$.

\begin{proposition}[Existence of a finite analytic refinement]
\label{prop:finite-refinement}
Fix a finite positive spectrum $\{\lambda_j\}_{j=1}^{r}$, an upper-tail query $q>\E Q$, and $\eps>0$. There exist a finite interval endpoint $b>q$, a finite number of cosine modes $N$, and a finite packet partition (allowing singleton packets) for which the analytic error bound in \eqref{eq:combined} is below $\eps$.
\end{proposition}

\begin{proof}
First choose $b$ sufficiently large that \eqref{eq:eps-support} is below $\eps/3$. This is possible because the tilted law has a moment generating function in a nontrivial right neighborhood of zero and hence its Chernoff bound tends to zero exponentially as $b\to\infty$. Hold this $b$ fixed. By \eqref{eq:phi-modulus}, $|\phi_{\theta}(t)|\to0$ polynomially as $t\to\infty$, while $\beta(t)\to r/2>0$. Therefore \eqref{eq:eps-cos} tends to zero as $N\to\infty$, so a finite $N$ makes it below $\eps/3$. Finally, split every packet to singletons if necessary. A singleton has $\delta_j=0$, so its packet remainder is exactly zero. Hence a finite partition makes $\eps_{\mathrm{pkt}}<\eps/3$. Summing the three terms proves the claim.
\end{proof}

\begin{remark}\label{claim:final-26-14}
Proposition \ref{prop:finite-refinement} proves existence of a finite representation meeting any analytic tolerance. It does not by itself prove that every possible greedy refinement schedule terminates optimally. The prototype planner refines the largest currently certified obstruction and is evaluated empirically below.
\end{remark}

\subsection{Adaptive K-COS algorithm}

The prototype accepts $(\lambda_1,\ldots,\lambda_r,q,\tau)$, where $\tau$ is a requested relative tolerance.

\begin{enumerate}[leftmargin=2em]
\item Solve $\kappa'(\theta_*)=q$.
\item Initialize an upper interval $b=q+L\sqrt{\kappa''(\theta_*)+\sqrt{\kappa^{(4)}(\theta_*)}}$ and a mode count $N$.
\item Retain the largest spectral spikes exactly and place the remaining sorted eigenvalues into coarse packets.
\item Evaluate $\rho_P(\theta_*+\ii u_k)$ and the packet remainder bound on the actual COS grid. Split packets until the propagated packet contribution is within its budget.
\item Evaluate the support/alias certificate from \eqref{eq:eps-support}--\eqref{eq:eta}.
\item Evaluate the omitted-mode certificate \eqref{eq:eps-cos}.
\item If the combined relative certificate exceeds $\tau$, refine whichever analytic error source dominates: split packets, increase $b$, or increase $N$.
\item Return the tail probability, the compiled spectral passport, and the error ledger.
\end{enumerate}

The resulting object is query dependent. Schematically,
\begin{equation}
(\Lambda,q,\tau)
\longmapsto
\mathcal P_{\Lambda,q,\tau}
\longmapsto
(\widehat p,\text{certificate}),
\label{eq:compiler-map}
\end{equation}
where $\mathcal P_{\Lambda,q,\tau}$ stores exact spikes, packet centers, centered moments, packet radii, and the selected COS parameters.

\subsection{Archived packet-prototype experiments}
These tables and three associated figures are the original September 2026
packet-prototype measurements.  Their returned analytic bounds exclude
floating-point roundoff.  The original prototype predates the interval
backend and independent inversion checks already presented above; the
following measurements are preserved rather than recast as fresh benchmarks.

\subsubsection{Exactly solvable chi-square stress test}

To isolate the effect of exponential tilting, we first used
\[
Q=0.02\chi^2_{80},
\]
whose survival probabilities are available directly from the chi-square distribution. Table \ref{tab:chi} compares a global COS calculation with query-centered tilted COS.

\begin{table}[htbp]
\centering
\caption{Relative error in an exactly solvable chi-square family.}
\label{tab:chi}
\begin{tabular}{rrrr}
\toprule
True tail $p$ & $N$ & Global COS rel. error & Tilted COS rel. error\\
\midrule
$10^{-6}$  & 64  & $3.56\times10^{-8}$ & $2.71\times10^{-14}$\\
$10^{-10}$ & 64  & $1.64\times10^{-4}$ & $6.59\times10^{-15}$\\
$10^{-10}$ & 128 & $2.03\times10^{-6}$ & $9.05\times10^{-15}$\\
$10^{-14}$ & 64  & $1.43$              & $3.53\times10^{-14}$\\
$10^{-14}$ & 128 & $6.15\times10^{-2}$ & $3.87\times10^{-14}$\\
\bottomrule
\end{tabular}
\end{table}

The test demonstrates the numerical motivation for tilting: evaluating a small tail directly under the shifted law is substantially more stable than reconstructing a global CDF and subtracting from one. This observation is not a novelty claim for exponential tilting itself.

\subsubsection{Genotype-derived spectra}

The second experiment used the existing Komorebi SKAT benchmark generator. Synthetic genotype matrices were centered and scaled, residualized against three covariates, weighted by $(1-\mathrm{MAF})^{24}$, and converted to the positive quadratic-form spectrum
\[
\lambda=\tfrac12\operatorname{eig}(AA^{\mathsf T}).
\]
Three independent $n=400$, $m=1200$ genotype sets yielded numerical rank $397$. Four upper-tail queries per spectrum were generated by choosing saddlepoint locations at fractions $0.70,0.80,0.85,0.90$ of the transform singularity. The comparison reference used exact-spectrum tilted COS with $N=2048$ and a large support parameter $L=20$. The adaptive method was asked for a relative analytic tolerance of $10^{-8}$.

Across the twelve queries:
\begin{itemize}
\item every requested $10^{-8}$ analytic certificate passed;
\item the worst observed relative error was $6.148608\times10^{-11}$;
\item the largest returned analytic relative bound was $7.917203\times10^{-9}$;
\item the planner selected $N\in\{64,128\}$, $L\in\{10,11,12,13\}$, and 11--16 packets.
\end{itemize}

Figure \ref{fig:coverage} plots the observed errors against the returned analytic bounds. Every point lies above the equality line because the vertical coordinate is the bound.

\begin{figure}[htbp]
\centering
\includegraphics[width=0.78\linewidth]{figures_packets/certificate_coverage.png}
\caption{Observed relative error versus the analytic certificate on twelve genotype-derived tail queries. The dashed line is equality.}
\label{fig:coverage}
\end{figure}

\subsubsection{Tolerance ladder}

At a reference tail probability near $8.99\times10^{-15}$, the same engine was requested to meet three tolerances. Table \ref{tab:ladder} shows how the representation refined automatically.

\begin{table}[htbp]
\centering
\caption{Adaptive refinement at an extreme-tail query.}
\label{tab:ladder}
\begin{tabular}{rrrrrr}
\toprule
Requested $\tau$ & $N$ & $L$ & Packets & Analytic bound & Observed error\\
\midrule
$10^{-6}$  & 128 & 10 & 10 & $3.56\times10^{-7}$  & $3.63\times10^{-12}$\\
$10^{-8}$  & 128 & 12 & 15 & $4.87\times10^{-9}$  & $9.14\times10^{-13}$\\
$10^{-10}$ & 256 & 14 & 24 & $8.21\times10^{-11}$ & $5.65\times10^{-14}$\\
\bottomrule
\end{tabular}
\end{table}

\begin{figure}[htbp]
\centering
\includegraphics[width=0.78\linewidth]{figures_packets/tolerance_ladder.png}
\caption{Requested tolerance, returned certificate, and observed error for one extreme-tail query.}
\label{fig:ladder}
\end{figure}

\subsubsection{Prototype runtime scaling}

Table \ref{tab:runtime} reports median CPU timings from the Python prototype. ``Adaptive total'' includes planning and analytic error checks. ``Compiled evaluation'' measures reuse of an already compiled spectral passport at the selected grid.

\begin{table}[htbp]
\centering
\caption{Prototype runtime scaling on genotype-derived spectra.}
\label{tab:runtime}
\begin{tabular}{rrrrrr}
\toprule
Rank & Packets & Exact tilted & Adaptive total & Compiled eval. & Compiled speedup\\
\midrule
397  & 13 & 6.66 ms & 10.45 ms & 0.549 ms & $12.1\times$\\
797  & 16 & 13.32 ms & 11.49 ms & 0.604 ms & $22.0\times$\\
1197 & 17 & 21.52 ms & 11.24 ms & 0.651 ms & $33.1\times$\\
\bottomrule
\end{tabular}
\end{table}

\begin{figure}[htbp]
\centering
\includegraphics[width=0.78\linewidth]{figures_packets/runtime_scaling.png}
\caption{Median runtime of exact tilted COS, adaptive K-COS with analytic ledger including planning, and evaluation of an already compiled passport.}
\label{fig:runtime}
\end{figure}

The planning overhead matters. At rank 397, one adaptive analytic-ledger query is slower than one exact tilted-COS query. The crossover appears between ranks 397 and 797 in this prototype, while compiled reuse is much faster. No universal speedup claim is therefore made.


\section{Memory planning law}\label{sec:final-27}

Let \(d\) be the number of bytes per floating scalar, \(c\) the resident
variant count, \(s\) the Krylov width, and \(p\) the covariate rank.
Ignoring library workspace and allocator overhead, a streamed sample-space
action has working memory on the order of
\[
\boxed{
d\,[nc+2ns+cs+np+O(ps+c)].
}
\]

The full factor costs
\[
dnm.
\]

For fixed \(n,p,s,c\),
\[
\boxed{
\frac{M_{\rm stream}}{M_{\rm full}}
=
O\!\left(
\frac{c}{m}
+
\frac{s}{m}
+
\frac{p}{m}
\right),
}
\]
so device memory remains bounded as total variant count grows.

\begin{table}[H]
\centering
\caption{Algebraic float64 memory model for \(n=2000\), \(p=4\),
Krylov width \(s=16\), resident variant block \(c=1024\).}
\begin{tabular}{rrrr}
\toprule
\(m\) & full factor (MiB) & streamed \(KV\) (MiB) & tiled Gram (MiB)\\
\midrule
10,000 & 152.6 & 16.308 & 8.555\\
50,000 & 762.9 & 16.308 & 8.555\\
100,000 & 1,525.9 & 16.308 & 8.555\\
500,000 & 7,629.4 & 16.308 & 8.555\\
1,000,000 & 15,258.8 & 16.308 & 8.555\\
\bottomrule
\end{tabular}
\end{table}


\begin{figure}[H]
\centering
\includegraphics[width=.84\linewidth]{figure_03_memory_law.png}
\caption{The memory law. The complete factor grows linearly with total variant count $m$, whereas the streamed kernel action and the two-dimensional tiled construction remain bounded when resident block size, Krylov width, and covariate rank are fixed.}
\label{fig:memory-law}
\end{figure}

\section{Numerical verification}\label{sec:final-28}

All identities were compared against explicitly formed reference matrices.

\begin{table}[H]
\centering
\caption{Partition-invariance checks.}
\begin{tabular}{lrrr}
\toprule
partition & blocks & Gram relative error & action relative error\\
\midrule
coarse 4 & 4 & \(2.91\times10^{-16}\) & \(5.99\times10^{-16}\)\\
fine 64 & 32 & \(3.39\times10^{-16}\) & \(6.00\times10^{-16}\)\\
irregular & 9 & \(3.34\times10^{-16}\) & \(5.87\times10^{-16}\)\\
singleton-front & 21 & \(2.43\times10^{-16}\) & \(5.82\times10^{-16}\)\\
permuted, block 73 & 29 & \(3.52\times10^{-16}\) & \(5.89\times10^{-16}\)\\
\bottomrule
\end{tabular}
\end{table}

A hierarchical coarse-to-fine refinement gave relative discrepancy
\[
3.17\times10^{-16}.
\]

For an explicitly PSD polyphase operator
\[
H=L^\top L
\]
with block-bidiagonal \(L\), the smallest eigenvalue was approximately
\[
2.86,
\]
and the largest norm observed outside the exact \(k\)-step light cone was
\[
\boxed{0}.
\]

The QR-core update reproduced the full singular spectrum with relative
\(\ell_2\) discrepancy
\[
5.89\times10^{-16}.
\]


\section{Method synthesis: composable decompositions}\label{sec:final-29}

The solver is organized around four independent decomposition axes:
\[
\boxed{
\mathcal D
=
\mathcal D_{\rm storage}
\circ
\mathcal D_{\rm phase}
\circ
\mathcal D_{\rm rank}
\circ
\mathcal D_{\rm spectrum}.
}
\]
This display describes the pipeline schematically, not four endomorphisms
proved to commute pairwise.  The precise commutation theorem below is for
exact storage refinement and a fixed structural phase partition.  Full range
capture or residual completion preserves the nonzero spectrum exactly; a
truncated packet representation instead preserves the requested tail within
its explicit analytic error bound.

They answer different questions.  Storage decomposition controls what is
resident.  Phase decomposition controls where information can propagate.
Rank decomposition controls how many directions must be retained.  Spectral
packetization controls the cost of the final deep-tail transform.  The central
structural fact is that these representations can be changed independently
without changing the exact operator when their respective invariance
conditions are satisfied.

\subsection{Storage-phase commutation}

Fix a structural phase partition
\[
\{1,\ldots,m\}=\mathcal P_1\sqcup\cdots\sqcup\mathcal P_H
\]
and let each phase be subdivided into arbitrary storage blocks
\[
\mathcal P_h
=
J_{h1}\sqcup\cdots\sqcup J_{hB_h}.
\]
Let $S$ denote this physical storage subdivision.  The phase projectors
$P_h$ are defined by the structural phases, not by $S$.

\begin{theorem}[Storage-phase commutation theorem]\label{claim:final-29-1}
Let
\[
A=M_XGW^{1/2},
\qquad H=A^\top A,
\]
and let $S$ and $S'$ be any two exact storage refinements that respect the
same structural phase partition.  Then:
\begin{enumerate}[label=(\roman*)]
\item the exact operator $H$ is identical under $S$ and $S'$;
\item every phase coupling $P_hHP_g$ is identical under $S$ and $S'$;
\item the exact phase graph and every $k$-step polyphase light cone are
identical under $S$ and $S'$;
\item any exact matrix-free Krylov method applied through streamed products
produces the same Krylov subspaces in exact arithmetic under $S$ and $S'$.
\end{enumerate}
\end{theorem}

\begin{proof}
Storage refinement changes only the order and grouping of finite sums.
By partition refinement invariance, the streamed products reconstruct the
same $A$, $A^\top$, and therefore the same $H=A^\top A$.  Since each
$P_h$ is fixed by the structural phase partition, every block
$P_hHP_g$ is unchanged.  The phase graph is defined only by these blocks, so
its path structure and light cones are unchanged.  Finally, an exact
matrix-free Krylov method depends only on repeated applications of the same
operator to the same starting block; hence its generated subspaces are
identical.
\end{proof}

\begin{corollary}[Adaptive resident-block scheduler]\label{claim:final-29-2}
An implementation may change resident block size at runtime---including
recursively splitting a phase block after an allocation failure---without
changing the exact phase graph or exact Krylov trajectory, provided all
required storage blocks are processed exactly.
\end{corollary}

This separates \emph{mathematical locality} from \emph{physical storage
granularity}.  A phase may remain fixed as an LD window or annotation class
while its physical representation is subdivided according to free GPU
memory, scalar precision, transfer granularity, or stream count.

\subsection{Frontier-to-I/O theorem}

\begin{theorem}[Exact frontier-to-I/O reduction]\label{claim:final-29-3}
Suppose a validated exact phase graph of the fixed operator $H=A^\top A$
is already supplied and the initial Krylov block is supported on a phase set
$F_0$.  Keep covariate projections, weights, and structural phase metadata
fixed, and assume all factor-action blocks required for the retained phase
operator are available.  At depth $k$, let $F_k$ be the graph
$k$-neighborhood of $F_0$.  Then no storage block lying wholly in a phase
outside $F_k$ can affect the $k$-step Krylov vector.  Such a block need not be
loaded when evaluating the requested power through the restricted factor
$A_{F_k}$.  Storage blocks mixing active and inactive phases require a
separate refinement before this omission is valid.
\end{theorem}

\begin{proof}
For every phase $h\notin F_k$ and every starting phase $g\in F_0$, the
polyphase light-cone theorem gives $P_hH^kP_g=0$.  Storage-phase commutation
allows each active phase operator to be evaluated through arbitrary exact
storage blocks.  Every path of length at most $k$ beginning in $F_0$ stays in $F_k$.
Consequently restriction to $P_{F_k}HP_{F_k}=A_{F_k}^\top A_{F_k}$ preserves
the requested $k$-step power on the initial support, including its zero
outside output.  Blocks wholly outside the retained phase set may therefore
be omitted from this restricted factor evaluation.
\end{proof}

The theorem turns an algebraic support statement into an execution law:
\[
\boxed{
\text{Krylov depth}
\longrightarrow
\text{reachable phases}
\longrightarrow
\text{storage blocks that must be transferred}.
}
\]
Discovery or verification of the phase graph may itself require all inputs.
The theorem does not claim minimum total I/O, omission of mixed-phase blocks,
or correctness for an unverified approximate graph.  The scheduler is constrained
by the exact operator graph rather than solely by available memory.

\subsection{Full method stack}

The complete positive-form solver may therefore be written as
\[
\boxed{\begin{aligned}
A
&\longrightarrow \text{resident storage blocks}\\
&\longrightarrow \text{structural phase frontier}\\
&\longrightarrow \text{phase-local Krylov/range core}\\
&\longrightarrow \text{residual moment certificate or exact completion}\\
&\longrightarrow \{\lambda_j\}\\
&\longrightarrow \text{centered spectral packets}\\
&\longrightarrow \log P(Q\ge q).
\end{aligned}}
\]

Three invariances make the stack composable:
\begin{enumerate}
\item storage refinement does not change the represented operator;
\item structural phase scheduling does not depend on physical block size;
\item once the spectrum or a certified spectral enclosure is obtained, the
      tail engine is independent of the application that produced it.
\end{enumerate}

The resulting method can be aggressive during basis discovery and still return
to exactness.  Approximate phase pruning is controlled by query and operator
budgets; unresolved range structure is measured by trace, square-trace, and
spectral-cap invariants; and when the remaining rank is small, streamed
residual completion restores the exact nonzero spectrum before the final tail
evaluation.

\subsection{Method summary}

The resolved synthesis is
\[
\boxed{\begin{aligned}
&\text{exact streamed factor actions}
+\text{ structural light cones}
+\text{ certified pruning}\\
&+\text{ phase-local rank}
+\text{ residual moment certificates}
+\text{ exact residual completion}\\
&+\text{ spectral packets}
+\text{ log-domain extreme-tail evaluation}.
\end{aligned}}
\]

The central design principle is separation of representation from semantics:
changing how an operator is stored, scheduled, compressed, or evaluated must
not silently change the mathematical quadratic form.  Each approximation
layer therefore either carries an explicit certificate or terminates in an
exact completion step.


\section{Prior work and the contribution boundary}\label{sec:prior-work}
Characteristic-function methods for quadratic forms substantially predate
this engine.  Imhof's 1961 paper is a foundational historical comparator
\cite{Imhof1961}.  Davies' Algorithm AS~155 treats linear combinations of
noncentral chi-square variables with signed coefficients and an optional
Gaussian term, using characteristic-function inversion and explicit
truncation/integration controls \cite{Davies1980}.  Its documented difficult
low-degree cases are relevant to the retained low-rank limitations here.

The Fourier-cosine representation and the use of characteristic-function
values for cosine coefficients are due to Fang and Oosterlee
\cite{FangOosterlee2008}.  Their formulation distinguishes finite-domain
coefficients from the full characteristic-function coefficients used to
approximate them.  The support/folding argument above supplies the particular
positive-law alias bound needed here; a general COS accuracy theorem is not
inherited without checking its hypotheses.

No novelty is claimed for Gaussian spectral reduction, exponential tilting,
saddlepoint centering, Fourier-cosine inversion, Weyl bounds, QR updates, or
out-of-core multiplication themselves.  The proposed contribution is their
explicit integration with independently adjustable storage, structural,
range-residual, and query-specific centered spectral representations.  The
packet compiler in Section~\ref{sec:packets} is a concrete mechanism: its
local CGF remainder is propagated to the requested tail functional and
refined jointly with support and mode errors.  The finite refinement theorem
is an existence theorem, not a complexity guarantee or proof of planner
optimality.

The primary-source check for this final edition verified the listed COS
and Davies formulas and the Imhof bibliographic record.  It was not an
exhaustive priority search.  The archived independent Imhof-style inversion
checks support accuracy on their tested heterogeneous spectra; they do not
provide a universal runtime comparison with a mature Davies implementation.

\section{Guarantee levels and scope}\label{sec:final-31}

The framework contains several different evidence levels.  They are kept
separate because a proved algebraic identity, a finite numerical validation,
and a directed-rounding probability enclosure are different kinds of claims.

\begin{center}
\small
\begin{tabular}{@{}p{0.30\linewidth}p{0.20\linewidth}p{0.41\linewidth}@{}}
\toprule
Statement & Status & Meaning\\
\midrule
Block factor and streamed-action identities &
proved &
Exact finite-dimensional matrix identities; storage blocking introduces no
mathematical approximation.\\

Storage-phase commutation and exact light cones &
proved &
Physical block refinement does not change the structural phase graph; phases
outside the exact finite-depth frontier contribute zero.\\

Approximate light-cone and pruning certificates &
proved + tested &
Weak discarded couplings carry explicit finite-depth transfer and
operator-norm budgets.\\

Range residual and exact completion &
proved + tested &
The unresolved factor is measured directly; small residual rank can be
completed exactly.\\

Positive exact-spectrum tail enclosure &
directed rounding + proved remainders &
For rank \(>2\), the finite transform is enclosed numerically and enlarged by
proved support/alias and infinite-mode bounds.\\

Accuracy-driven positive tail solver &
certified / fail-closed &
A result is accepted only after
\(p_+/p_- - 1\le\varepsilon_{\rm requested}\) is proved.\\

Positive noncentral engine &
exact reduction + tested tail solver &
The Gaussian form is reduced to a deterministic shift plus positive
noncentral chi-square components.\\

Signed / indefinite engine &
proved reduction + analytic-ledger tests &
High-rank and finite-endpoint cases passed analytic error checks; low-rank mixed-sign forms
remain limited by slow Fourier decay of the present COS certificate.\\

Binary64 factor to rigorous spectrum &
implemented at moderate dimension &
Candidate eigensystems are post-certified by directed-rounding orthogonality
and reconstruction bounds, then propagated monotonically into the tail.  The
remaining issue is scalable streamed/subspace verification rather than the
basic correctness theorem.\\
\bottomrule
\end{tabular}
\end{center}

Three boundaries remain material:

\begin{enumerate}
\item scale the verified matrix-to-spectrum bridge to streamed range-core and
residual-completion linear algebra without materializing a full dense basis;
\item replace the conservative low-rank mixed-sign Fourier-tail certificate
with a rigorous convolution or quadrature route;
\item implement the proven streamed and residual-completion algorithms in the
final native GPU runtime and measure transfer/compute crossover on hardware.
\end{enumerate}

These are boundaries of the implementation, not hidden assumptions in the
certified positive matrix-to-tail result at moderate dimension.


\section{Conclusion}\label{sec:final-32}

The quadratic-form problem can be separated into representation, structure,
spectral recovery, and tail evaluation.

At the representation level, factors can be streamed in blocks without
changing the operator.  At the structural level, phase graphs constrain
finite-depth propagation and allow certified scheduling.  At the spectral
level, unresolved range structure is measured rather than guessed and can be
completed exactly when its rank becomes small.  At the probability level,
exponential tilting and spectral packets make extreme tails computationally
tractable, while directed-rounding interval arithmetic and analytic remainder
bounds turn the positive exact-spectrum calculation into an explicit
probability enclosure.

The resulting numerical contract is deliberately stronger than
``approximately evaluate a small probability.''  For the rigorous positive
backend, the user supplies an accuracy requirement.  Internal resolution is
refined until that requirement is proved, or the solver returns an
uncertified interval rather than claiming success.

The deepest archived September 2026 exact control reaches
\[
p=10^{-5000}.
\]
At that point the archived accuracy-driven run reported relative multiplicative
uncertainty below \(10^{-18}\).  This establishes that extreme-tail depth by
itself is not the limiting numerical variable; the meaningful limit is whether
the complete error ledger can be made smaller than the requested tolerance.

The matrix-to-spectrum bridge now closes that upstream correctness gap for
moderate binary64 factors by post-certifying an otherwise untrusted candidate
eigensystem.  The remaining rigorous frontier is primarily one of scalable
representation: carry the same enclosure through streamed range cores and
residual completion without forming a full dense eigenbasis, while extending
the fail-closed accuracy contract to difficult low-rank mixed-sign forms.



\begin{thebibliography}{9}
\bibitem{Imhof1961} J.~P. Imhof, Computing the distribution of quadratic forms
in normal variables, \emph{Biometrika} 48(3--4) (1961), 419--426.
\href{https://doi.org/10.1093/biomet/48.3-4.419}{doi:10.1093/biomet/48.3-4.419}.
\bibitem{Davies1980} R.~B. Davies, Algorithm AS~155: The distribution of a
linear combination of $\chi^2$ random variables,
\emph{Applied Statistics} 29(3) (1980), 323--333.
\href{https://doi.org/10.2307/2346911}{doi:10.2307/2346911}.
\bibitem{FangOosterlee2008} F. Fang and C.~W. Oosterlee, A novel pricing
method for European options based on Fourier-cosine series expansions,
\emph{SIAM Journal on Scientific Computing} 31(2) (2008), 826--848.
\href{https://doi.org/10.1137/080718061}{doi:10.1137/080718061}.
\end{thebibliography}

\appendix

\section{Reference pseudocode}\label{sec:final-A}

\begin{lstlisting}[language=Python]
def stream_KV(block_source, V, project=lambda x: x):
    U = project(V)
    Y = zeros_like(U)
    for A_block in block_source:
        Y += A_block @ (A_block.T @ U)
    return project(Y)

def certified_tail(q, spectrum, relative_accuracy):
    # Refines support, COS modes, and interval precision.
    # Returns certified=False if the requested accuracy is not proved.
    return certified_positive_logsf(
        q, spectrum, relative_accuracy=relative_accuracy
    )
\end{lstlisting}

\section{Reproducibility and evidence package}\label{sec:final-B}

The release accompanying this manuscript contains the executable engine,
interval backend, deterministic regression tests, numerical evidence tables,
and proof-audit notes used to support the quantitative claims.

The principal production-facing files are:

\begin{itemize}
\item \texttt{quadratic\_form\_engine.py}: general central/noncentral interface;
\item \texttt{accuracy\_driven.py}: fail-closed tolerance-driven controller;
\item \texttt{interval\_kcos.py}: directed-rounding positive-tail backend;
\item \texttt{kcos\_v2\_hardened.py}: log-domain adaptive positive-tail engine;
\item \texttt{generalized\_tail.py}: signed and noncentral generalized forms;
\item \texttt{residual\_completion.py}: exact completion of unresolved factor
range.
\end{itemize}

Evidence tables record exact-control, realistic-spectrum, signed-form,
noncentral, interval-enclosure, and requested-accuracy experiments.  The
release also records negative results and corrected implementation defects so
that failed reference calculations or superseded simulation defects are not
silently reused as evidence.



\section{Fresh final-edition validation}\label{app:fresh-validation}
\input{final_edition_validation.tex}

\section{Lean claim-coverage contract}\label{app:lean-contract}
\input{lean_coverage.tex}

\section{Source and figure provenance}\label{app:provenance}
The traditional text incorporates the full
\texttt{FINAL\_MANUSCRIPT.tex} from
\path{Komorebi_Quadratic_Form_Engine_v1.9_Verified_Matrix_to_Tail_Source}
and the mathematical and experimental core of the standalone
\path{K-COS_Query_Adaptive_Spectral_Packetization_v0.1_Source}
archive. Original matrix-engine figures remain unchanged under
\path{../engine/figures_final/}; the three original packet-prototype figures
are copied unchanged under \texttt{figures\_packets/}. The final
\texttt{figures\_final/} directory contains 24 layout derivatives and four
unchanged originals. Derivatives move annotations and legends into reserved
space, wrap labels, distinguish coincident series and separate the two
metrics in Figure~\ref{fig:dual-budget-tail}. Their numerical inputs
are the archived tables; these layout changes add no empirical evidence.

The raw plotted vectors for Figures~\ref{fig:phase-local-rank},
\ref{fig:qr-core-equivalence} and \ref{fig:spectral-packets} were not
archived. Their replacement visuals use, respectively, the exact recorded
phase-rank audit, the exact recorded QR discrepancy audit, and the original
rounded packet-center and second-moment labels with the formal packet
definitions. No hidden singular-value or eigenvalue samples are inferred.
The portable scripts \texttt{relayout\_figures.py} and
\texttt{relayout\_late.py} record source hashes, plotted coordinates and
preservation checks in their figure-layout reports. Underlying tables and
original illustrations remain in the executable source and historical
archives accompanying the combined release.

The final-edition modifications add the scope contract, integrate packet
proofs, justify the integrated cosine identity by folding and Fej\'er means,
make QR complement orthogonality explicit, include resident covariate-tile
storage, clarify graph-knowledge assumptions, and distinguish analytic
ledgers from directed numerical enclosures.  Exact source hashes and the
fresh validation manifest are provided with the release.

\end{document}
